\section{Computations}\label{sec:verification}

Every statement of the paper is proved in the text, and no proof rests on
the items below. They repeat the finite arithmetic of the proofs in exact
arithmetic: rational, integer or exterior-algebra computations, and exact
linear algebra over number fields and over function fields in one variable,
or, for the integrals of theta functions of part (H) of item (LXXII), in ball
arithmetic, which returns for every number an interval guaranteed to
contain it. Some items also record data beyond what the paper uses, such as
the searches of items (XXXII), (LVII), (LVIII) and (LXXVI) and the ranks of
item (LXXII); the paper does not cite them. The quadrature of item (LXXII)
and the random tests in items (LXXI) and (LXXIII) are in floating point. A rank computed modulo a prime is used only
as a lower bound for the rank over $\QQ$, since reduction can only lower a
rank, and the bound is matched by an exact computation in the other
direction. Items (XXIX), (XXXI) and (XL) are $\Ext$
computations over a polynomial ring. The programs that carry out the
computations are in the directory \texttt{code/} of this archive, and the driver
\texttt{verify\_all.py} runs them item by item. For each item we describe what is
computed, how it is computed, and which statements of the paper it supports.

\begin{enumerate}[label=\textup{(\Roman*)},leftmargin=*]
\item \emph{The Fourier--Mukai transform on powers of the polarisation.}
  The exterior algebra $\bigwedge^{*}(V\oplus V^{\vee})$ of
  \Cref{not:model} is constructed explicitly for $g\le4$, with a basis
  indexed by subsets of the generators and structure constants given by the
  shuffle sign. The class $\ell$ of \eqref{eq:ell}, the forms $\theta$ and
  $\theta'$ of \eqref{eq:thetamodel} and \eqref{eq:thetapmodel}, and the
  exponential $e^{\ell}$ are formed, and $\Phi_{\cP}(\theta'^{\,k})$ is
  computed by extracting the coefficient of the top form in the dual
  variables. The result is an exact rational multiple of $\theta^{\,g-k}$
  with no other component, and the multiple agrees with \eqref{eq:fmtheta} in
  every case.

\item \emph{The sign.} The proof of \Cref{thm:fmtheta} uses that the sign
  $\sigma(g,k)$ is the same for all index sets $T$. This is checked as a
  permutation computation for every $g\le12$ and every $0\le k\le g$,
  together with the closed form \eqref{eq:sign}; all values agree.

\item \emph{The character separation.} The grading \eqref{eq:grading} is
  formed for $n\le4$ and the positions of the Weil line and of $\eta^{n}$ are
  computed, confirming \Cref{lem:NSchar}, \Cref{lem:weilchar} and the
  disjointness of \Cref{thm:character}.

\item \emph{The numerical comparison.} The values of $2n^{2}+2n$, $2^{2n}$
  and $3^{2n}$ are tabulated for $2\le n\le8$, confirming
  \Cref{lem:chicompare}. As counterexamples to the inference discussed in
  \Cref{rem:retract}, the item also records the polarisation types $(1,3)$
  and $(1,5)$ on an abelian surface, with Euler characteristics $3$ and $5$.

\item \emph{The semiregularity target.} The identity
  $\sum_{q}\binom{2n}{q}\binom{2n}{q+2}=\binom{4n}{2n-2}$ of
  \Cref{prop:target} is verified for $2\le n\le7$, with the growth rates of
  \Cref{cor:countnotbind}.

\item \emph{The coordinate model.} The model of \Cref{setup:model} is built
  over $\QQ(\sqrt{-d})$ for the pairs $(n,d)$ of a table covering $n\le3$ and
  $d\in\{1,2,3,5,7\}$, and eight statements are checked for each: that the
  $u^{s}_{j}$ of \eqref{eq:eigenvec} are eigenvectors of $S$; that $\omega_{+}+\omega_{-}$ and
  $(\omega_{+}-\omega_{-})/\delta$ are rational and agree with the closed form
  \eqref{eq:omegagens}; that $\iota(\tau)^{*}$ multiplies $\omega_{\pm}$ by
  $\tau^{2n}$ and $\bbar{\tau}^{\,2n}$ at five values of $\tau$; that
  $\omega_{\pm}$ is of Hodge type $(n,n)$ for the balanced structure and of
  type $(p,2n-p)$ for every other one, which is \Cref{prop:signatureforced};
  that $\eta$ is rational, of type $(1,1)$ and a norm eigenclass; that the two
  Riemann relations hold with the sign vector of \Cref{setup:model}; that
  $\eta^{2n}\ne0$; and that the supports of $\omega_{1},\omega_{2}$ and
  $\eta^{n}$ are disjoint, which is \Cref{prop:disjoint}.

\item \emph{The Hodge class count.} For $n\le4$ and every $(r,s)$, the
  dimension of the space of $\mathfrak{sl}_{2n}$-invariants in
  $\bigwedge^{r}V_{+}\otimes\bigwedge^{s}V_{+}^{\vee}$ is computed in exact
  rational arithmetic, as the kernel of the raising operators on the
  weight-zero subspace. This gives $\dim\Hdg^{n}=3$ and $\dim\Hdg^{k}=1$ for
  $k\ne n$, which is \Cref{thm:hdgcount}, and the codimension count of
  \Cref{prop:weildomain}.

\item \emph{The Clifford algebra.} The algebra $C(T,q)$ is built for forms
  of signature $(3,b-3)$ with $b\le8$. The computation checks the Clifford
  relations and $v^{2}=q(v)$ for a general $v$; that $J^{2}=-1$ and that
  left multiplication by $J$ squares to $-\id$ on $C^{+}$, which is
  \Cref{lem:cliffJ}; and that $C^{-}C^{+}C^{-}\subseteq C^{+}$, which is why
  the map $\Psi$ of \Cref{prop:ksembed} takes values in $\End(C^{+})$. It
  also tabulates the Hodge numbers of \Cref{prop:h2pieces}, together with the
  solution $n=1$ of \Cref{prop:normpart}.

\item \emph{The family and the secant plane.} The computations verify
  \Cref{thm:explicitweil} symbolically for $n\le5$ and several $d$, and
  derive the closed forms of \Cref{thm:plane} and
  \Cref{thm:secantinvariants}.

\item \emph{The split member.} The computation builds
  $A=X\times\widehat{X}$ with the action \eqref{eq:Mdef} in coordinates and
  makes three checks, the first two independent of each other. First, exact
  rational Gaussian elimination expresses each of $\omega_{1},\omega_{2}$ as
  a rational combination of the degree $2n$ monomials in
  $\beta,\widehat{\beta},\ell$, and the solution is recorded; this succeeds
  for $n\le3$ and $d\in\{1,2,3,7\}$. Second, \Cref{thm:splitclosed} is
  verified directly: $\gamma\mp\sqrt{-d}\,\ell$ is an eigenclass for
  $\tau^{2}$, respectively $\bbar{\tau}^{\,2}$, and \eqref{eq:splitclosed}
  holds exactly, for $n\le5$ and several $d$. Third, the formulas of
  \Cref{ex:splitsmall} are checked as stated, as equalities with the
  integral generators of \Cref{thm:explicitgens} for the $K$-basis
  $x_{1},\dots,x_{2n}$, at $n=2$ for $d\le3$ and at $n=3$ for
  $d\in\{1,2,3,7\}$.

\item \emph{The reach of the base case.} The computation checks the four
  assertions of \Cref{prop:splitgeom}. It verifies $M^{2}=-d$ and
  $M^{\top}Q+QM=0$ entry by entry for the block form
  $Q=\beta\oplus(\widehat\beta/d)$ on $H_{1}(A,\QQ)$, for $n\le4$, and checks,
  for six pairs $(n,d)$, that the action induced on $H^{2}(A,\QQ)$ multiplies
  $d\beta+\widehat\beta$ by $d$, multiplies $\gamma$ and $\ell$ by $-d$, and
  has $\beta+d\widehat\beta$ as an eigenclass only when $d=1$. It computes
  $\det H=(-d)^{n}(d_{1}\cdots d_{n})^{2}$ for ten choices of $(n,d,\beta)$,
  including non-principal $\beta$, and certifies in each case that
  $(-1)^{n}\det H$ is a norm from $K$, together with the two facts the proof
  uses, that $\det B$ is a square and that $d$ is a norm. It checks the
  codimension formula $n^{2}-n(n+1)/2=n(n-1)/2$ for $n\le10$. Finally it
  computes both the dimension of the space of $\mathrm{Sp}$-invariants and
  the rank of the set of monomials $\beta^{i}\widehat{\beta}^{j}\ell^{k}$ in
  degree $2n$, finding rank three in degree two for $n\le4$ and agreement in
  the middle degree for $n\le3$.

\item \emph{Quaternionic multiplication.} The computation builds the
  quaternion algebra $B=(-d,b)$ and the module $B^{n}$ in coordinates and
  verifies \Cref{thm:quatweil} for nine choices of $(n,d,b,a)$ with $n\le3$:
  the
  relations $i^{2}=-d$, $j^{2}=b$, $ij=-ji$; that $E$ of \eqref{eq:quatE} is
  alternating and nondegenerate with $i$ anti-self-adjoint and $j$
  self-adjoint; that the hermitian form on the $K$-basis is
  $\mathrm{diag}(2a_{k}d,-2a_{k}bd)$ with signature $(n,n)$; the determinant
  formula and hence the class $b^{n}$; and that $g(x,y)=E(jx,y)$ is
  alternating, $K$-bilinear and nondegenerate, which is the hypothesis of
  \Cref{thm:divcriterion}. It then checks \Cref{thm:everyfamily} numerically:
  for each of several fields, every nontrivial class in the tested range is
  realised in every dimension $n\le8$ with signature $(n,n)$, by a single
  factor when $n$ is odd and by the product of \Cref{lem:products} when $n$ is
  even. Finally it records the two dimensions behind \Cref{rem:ntwo}: the
  quaternionic locus has dimension $n(n+1)/2$ inside a period domain of
  dimension $n^{2}$, so its codimension is $n(n-1)/2$, which is checked for
  $n\le10$ and is $1$ exactly at $n=2$.

\item \emph{The semiregularity map of a sum of line bundles.} The
  computation builds the split member with $X=E_{i}^{n}$, so that the Hodge
  decomposition is defined over $\QQ(i)$ and every step is exact, and
  verifies the following. The six degree four monomials in
  $\beta,\widehat\beta,\ell$ are independent for $n\le4$, which is
  \Cref{lem:sixmonomials}. For six pairs $(n,d)$, $\theta^{+}\theta^{-}$ is
  never a multiple of $\eta^{2}$, which is \Cref{lem:tensorrank}. The
  coefficient of $\theta^{+}\theta^{-}$ in $\ch_{2}$ is
  $\sum_{i}m_{i}\Nm(u_{i})$, the identity behind \Cref{thm:positivity}, and
  for five pairs $(n,d)$ no multiset with positive multiplicities satisfies
  the resulting equations with some $u_{i}$ nonzero. The map $\Phi$ of
  \Cref{lem:splitsemireg}(iii) is injective at $n=2$ for $s\le3$ and has rank
  $22$ rather than $24$ at $s=4$. An exhaustive search over the classes
  $a\beta+b\widehat\beta+c\ell$ with $|a|,|b|,|c|\le3$ and $d\in\{1,2\}$
  finds no signed tuple of at most three line bundles with a flat Chern
  character and a nonzero Weil component, and every signed tuple of at most
  four in the same box with these two properties has exactly four terms and a
  map $\Phi$ with nonzero kernel, of dimension $2$, $3$, $4$ or $6$ at $d=1$
  and $3$ or $6$ at $d=2$. At $n=3$ the computation finds the
  Prouhet--Tarry--Escott pair of \Cref{lem:pte} on the norm circle $\Nm=4$ of
  $\QQ(\sqrt{-3})$ by exhaustive search, and verifies that the six line
  bundles of \Cref{prop:explicitobject} have $\ch_{k}=0$ for $k\ne3$ and
  $\ch_{3}=12\,\omega_{1}$ and that every difference of the six classes is
  nondegenerate of index three. A second computation repeats the rank
  computations modulo a prime congruent to $1$ modulo $4$, which reaches
  $n=3$: it reproduces the $n=2$ ranks exactly, finds $\Phi$ injective at
  $n=3$ for a general choice of up to ten classes, and finds rank $75$ of
  $90$ for the explicit object.

\item \emph{The annihilator and the rank-zero obstruction.} Modulo the same
  prime, the annihilator of each rational generator of the Weil line in
  $H^{0,2}$ has dimension $n^{2}$ at $n=2,3,4$, which is
  \Cref{lem:annihilator}. The computation splits the semiregularity map of
  the explicit object of \Cref{sec:construction} by parity and finds an odd
  kernel of dimension $9=n^{2}$ equal to the annihilator, an even kernel of
  dimension $6$, and a total of $15$, against $0$ for a generic triple, which
  is \Cref{cor:pteobstructed}. It then computes the semiregularity rank for
  every configuration $s_{i}\gamma+c_{i}\ell$ with $\sum s_{i}=\sum c_{i}=0$
  in a box of side $7$ whose virtual Chern character lies on the Weil line:
  there are eighteen such configurations, and the rank is $75$ of $90$ for
  all of them, which is \Cref{rem:familyuniform}.

\item \emph{The Weil class under products.} The computation verifies, in
  exact arithmetic over $\QQ(\sqrt{-d})$ for $d\in\{1,2,3,5,7\}$, the two
  identities \eqref{eq:weilmult} for the integral generators
  \eqref{eq:omegagens} and every pair $n_{1},n_{2}\le3$; the closed form
  \eqref{eq:omegaprod} for $n\le4$ and the $K$-valued form
  $\omega(A_{1})\omega(A_{2})=\omega(A_{1}\times A_{2})$ over the same
  range; the $k$-factor identity, with constant $1$, for six partitions
  with $k=3$ and $k=4$ and $d\le3$; that the Weil class lies in $H^{2n}$, so that at
  $n=1$ it is a divisor class, which is \Cref{lem:basecase}; and the
  dimension comparison of \Cref{rem:prodscope}, that the product locus has
  dimension $n$ inside a family of dimension $n^{2}$ and is smaller than the
  quaternionic locus for every $n\ge2$.

\item \emph{The annihilator as the tangent space.} The computation builds a
  polarised Hodge structure of $(K,-1,n)$-Weil type directly over the
  Gaussian rationals, with no reduction modulo a prime, and checks five
  statements at $n=2,3,4$. First, the construction is what it claims to be:
  the polarisation has exactly two nonzero blocks on $V_{+}\times V_{-}$,
  each a perfect pairing of $n$-dimensional spaces, and the two rational
  generators of the Weil line are real of type $(n,n)$. Second, the
  annihilator of each generator in $H^{0,2}$ has dimension exactly $n^{2}$,
  which is \Cref{lem:annihilator} again, in a model built independently of
  the one used in item (XIV). Third, the polarisation-preserving first-order
  deformations have dimension $n(2n+1)=\dim\cA_{2n}$, and those commuting
  with $K$ have dimension $n^{2}=\dim\cD_{n,n}$. Fourth, a polarised
  deformation contracts the whole Weil line to zero exactly on that
  $n^{2}$-dimensional subspace, and one generator already cuts it out, which
  is \Cref{prop:firstorder}. Fifth, the map of \Cref{thm:anntangent} is
  injective, independent of the basis of $V_{+}^{1,0}$ used to write it, and
  has image exactly the annihilator. The five statements are checked again
  for eight distinct choices of the $n$-plane $V_{+}^{1,0}$ at $n=2$ and at
  $n=3$, and hold in all sixteen.

\item \emph{The quaternionic locus as a Lagrangian Grassmannian.} The
  computation verifies, over $\QQ$ in exact arithmetic and for $2\le n\le5$,
  the three statements of \Cref{prop:lagrangian}. First, the form
  $\omega_{\psi}(x,y)=E(x,\psi y)$ attached to a Rosati-symmetric $\psi$ is
  alternating on $V_{+}$, which is the identity $E(\psi x,y)=E(x,\psi y)$
  read twice. Second, the graph of a symmetric form is isotropic and the
  graph of a form that is not symmetric is not, so the chart of
  $\mathrm{LG}(n,2n)$ by symmetric forms describes the locus exactly and is
  not merely contained in it. Third, the chart has dimension $n(n+1)/2$,
  computed as the rank of the symmetric basis rather than quoted, so that the
  codimension in $\cD_{n,n}$ is $n(n-1)/2$; at $n=4$ these are $10$ and $6$.
  The computation also checks that a random $n$-plane, completed to a basis
  of $V_{+}$, is Lagrangian for the form that is standard in that basis,
  which is the construction behind \Cref{rem:onlyrational}: over the reals
  the loci sweep out the whole family, and only the rationality of $\psi$
  prevents them from covering it.

\item \emph{The size forced on an object carrying the class.} In the
  exterior algebra over $\QQ$, for $1\le n\le4$, the computation verifies that
  $\alpha_{\pm}^{2}=0$, hence $\omega_{1}^{2}=2\alpha_{+}\alpha_{-}$, that
  $\int\omega_{1}^{2}\ne0$, and that $\chi(E,E)=(-1)^{n}N^{2}\int\omega^{2}$
  for every rank $r$ and every multiplicity $N$ in the ranges tested, the rank
  cancelling out of the identity. This is \Cref{prop:objectsize}, and with it
  the lower bound $\sum_{i}\dim\Ext^{i}(E,E)\ge N^{2}|\int\omega^{2}|$. The
  constant is not dimension-free, since it depends on the normalisation of
  $\omega$ against $\eta$, so no growth rate in $n$ should be read off it;
  the bound needs only that the Weil class is not isotropic.

\item \emph{Infinitesimal rigidity of divisor classes.} The computation
  builds the four blocks of $H^{1,1}$ and the tangent space of $\cD_{n,n}$
  over $\QQ(i)$ and verifies, for $n=2,3,4$, the three parts of
  \Cref{prop:rigidity}. First, the real span of the classes killed by
  contraction has dimension exactly one and is the polarisation; this is
  checked against the full real space of $(1,1)$ classes, of dimension
  $4n^{2}$, which is $16$ at $n=2$ and $64$ at $n=4$. Second, the kernel
  attached to a norm-character class with block $a$ has dimension
  $n(n-\operatorname{rank}a)$ for every rank; at $n=4$ the four values are
  $12$, $8$, $4$, $0$. Third, a class drawn at random has kernel zero, so it
  rigidifies the family to a point. The dimension $n^{2}$ of the tangent
  space is recomputed in the same computation rather than assumed.

\item \emph{The invariants of the required object.} For $1\le n\le8$, every
  squarefree $d\le11$, every rank $0\le a\le4$ and every twist $|b|\le4$, the
  Chern classes of the Chern character $au+bv$ are computed by Newton's
  identities; every one of them is integral and equal to
  $\gamma_{k}t^{k}/k!$ for the integer recursion of
  \Cref{prop:notforbidden}(ii), which proves integrality in every degree. The
  computation also confirms the closed form $\Delta=a^{2}d+b^{2}$ of
  \Cref{thm:secantinvariants} and its positivity: nothing numerical forbids
  the object the construction needs.

\item \emph{Existence of a sheaf with the required Chern character.} The
  computation behind \Cref{thm:secantexists} is carried out exactly over
  $\QQ$ for $1\le n\le10$. For every squarefree $d\le11$ and every rank and
  twist in a range, it solves the Vandermonde system that expresses the
  rational point $au_{t}+bv_{t}$ of the secant plane as a rational
  combination of the Chern characters $e^{jt}=\ch(\cO_{X}(jt))$, clears
  denominators to obtain the least positive multiple $M$ and the integer
  witness \eqref{eq:lbwitness}, and then verifies that witness independently
  of the solution, by recomputing every moment $\sum_{j}n_{j}j^{k}$ and
  comparing it with $Mc_{k}$. It checks that the rank $\sum_{j}n_{j}$ of the
  witness equals $Ma$, which is the positive rank hypothesis of
  \Cref{lem:posrank}, and it records the least $M$ in each case: $6$ at $n=4$
  with $d=3$ and $a=b=1$, at most $2016$ over the whole range at $n=8$, and
  never $1$ once $n\ge7$, which is \Cref{rem:themultiple}. Finally it expands
  $\theta^{k}/k!$ in the exterior algebra of a symplectic lattice and finds
  every coefficient integral; this is why the coordinates $c_{k}$ are
  integers, and hence why the problem is a lattice problem.

\item \emph{Semiregular split objects.} The computation settles
  \Cref{thm:pterange} in exact integer arithmetic. It enumerates the norms
  represented by $\cO_{K}$ below a bound, forms for every set of $n+1$ of them
  the forced level sums \eqref{eq:levelsums} and the least integral scale,
  and tests them against the supply $r_{K}(N_{j})$ of elements of that norm;
  at $n+2$ norms the solution space is two-dimensional, and its integer
  points in the same box are swept coordinate by coordinate. For $4\le n\le8$
  and $d\le19$ nothing survives, with norms up to $110$ at $n=4$ and up to
  $30$ at $n=8$. At $n=3$ and norms up to $200$ exactly eleven sets survive,
  all at $d=1$ or $d=2$, of the three shapes recorded in
  \Cref{thm:pterange}(iii). The computation then runs the whole system
  \eqref{eq:pteconditions} on those sets: for each it enumerates every
  assignment of signs to the elements of $\cO_{K}$ at those norms with the
  prescribed level sums, splitting the four levels into two halves, one
  tabulated and the other enumerated against the table, so that the
  enumeration is complete; it finds no solution for eight of the sets. A
  second computation settles the other three, $d=1$ at $(13,52,117,130)$,
  $(17,68,153,170)$ and $(25,50,100,125)$, whose largest level has sixteen
  elements: it matches three levels against the fourth through a linear hash
  of the ten off-diagonal coordinates, re-checks every match exactly, and
  finds no solution. As a control it also reruns two of the eight. The
  largest search covers $4.29\times10^{14}$ sign patterns.
\item \emph{The quaternionic Weil cycle.} The computation checks
  \Cref{prop:quatdivisible} on the lattice $\cO_{b}^{2}$. It confirms the
  ratios $-b/3$ and $4b^{2}/3$ and the value $1/12$ for $d\in\{1,3\}$, two
  weight vectors and fifteen values of $b$ up to $101$; the decomposition
  $E_{b}=E_{u}+bE_{v}$, $g_{b}=bg_{1}$, $\lambda_{b}=b\lambda_{1}$ with $g_{1}$
  and $\lambda_{1}$ integral; the equalities $Z_{b}=b^{2}Z_{1}$ and
  $Z'_{b}=b^{2}Z'_{1}$ in $\wedge^{4}\Lambda^{*}$; the growth
  $\int\eta_{b}^{4}=b^{2}\int\eta_{1}^{4}$; and the effect of replacing $j$ by
  $\alpha j$. It then computes the saturation of the Weil plane in
  $\wedge^{4}\Lambda^{*}$ and the sublattice spanned by products of two
  integral divisor classes, for $d\in\{1,2,3,7\}$, three weight vectors and
  seven values of $b$, and finds Gram matrix $\mathrm{diag}(8d,8d^{2})$ and
  index $2(a_{1}a_{2})^{2}$ throughout.
\item \emph{The divisor route on a fixed lattice.} The computation checks
  \Cref{prop:pencil}, which is proved by hand for every $d$, on
  $\Lambda=\cO^{2}$ with $j^{2}=1$, for $d\in\{1,3\}$. It computes the
  lattice of $K$-bilinear classes, of rank $12$, checks $t(x)=0$ on it and the
  identity $-\int xy\,\eta^{2}\big/\!\int\eta^{4}=\mathrm{tr}(\phi_{x}\phi_{y})/24$
  on all $78$ pairs of basis vectors, and finds signature $(8,4)$ for $I$. It
  builds the pencil $x_{k}=E(\phi_{k}\cdot,\cdot)$ from its definition and
  checks at eight values of $k$ that $x_{k}$ is integral, primitive and in
  $R$, that $\phi_{k}^{2}=(1+k^{2})/(4d^{2})$, that
  $I(x_{k})=\beta(k)\int\eta^{4}/3$, and that $\int\eta^{4}=384d^{4}$. For
  $0\le k\le10$ it saturates $\QQ\eta+\QQ x_{k}+\QQ x_{k}(i\cdot,\cdot)$ in
  $H^{2}(\Lambda,\ZZ)$, checks that the result is
  $\ZZ\tfrac{1}{2d}\eta\oplus\ZZ x_{k}\oplus\ZZ x_{k}(i\cdot,\cdot)$ and
  that $I$ is positive definite on it, and finds the minimum of $I$ off
  $\QQ\eta$, by an exhaustive search in a box derived from the inverse Gram
  matrix, to be $32d^{2}(1+k^{2})$. It forms the $2\times2$ minors of the
  matrix with rows $x_{k}$ and $x_{k}(i\cdot,\cdot)$ as polynomials in $k$,
  finds resultants of pairs of them with gcd $1$, so the minors are coprime at
  every integer $k$, and checks this directly for $|k|\le200$. Finally, for
  $0\le k\le4$ it computes the centraliser of $i$ and $\phi_{x_{k}}$ in
  $\mathfrak{sp}(V,E)$, of dimension $10$, and its invariants in
  $\bigwedge^{2}V^{*}$, of dimension $3$; and at $k=0$, where
  $\beta(0)=(2d)^{-2}$ is a rational square, it constructs a rational complex
  structure $J_{0}$ on the pencil member, commuting with $i$ and $\phi_{0}$
  and compatible with $E$, and checks that $[\mathfrak{c},J_{0}]$ has
  dimension $6$ and that $[\mathfrak{c},J_{0}]$ together with its brackets
  spans all of $\mathfrak{c}$, which is the Lie algebra step in the proof of
  \Cref{prop:pencil}(iv).
\item \emph{The obstruction map of a split object.} The computation works in
  the model of item (XIII), modulo the same prime, and checks
  \Cref{prop:splitobstruction}. For $(n,d)=(2,1),(2,3),(3,3)$ it computes the
  space $T$ of first-order deformations killing $\eta$ and both generators of
  the Weil line, of dimension $n^{2}$, and the space $T_{\mathrm{spl}}$ of
  deformations keeping $\beta,\widehat\beta,\ell$ of type $(1,1)$, of
  dimension $n(n+1)/2$, and checks that $T_{\mathrm{spl}}\subseteq T$. As a
  control, replacing $\eta$ by $\beta+d\widehat\beta$ makes $T$ collapse to
  dimension $n(n+1)/2$ for $d>1$, which is the mismatch of conventions
  discussed after \Cref{setup:construction}. For three classes spanning
  $\NS(A_{0})_{\QQ}$ together with $\eta$, the obstruction map
  $v\mapsto(v\lrcorner\lambda_{i})_{i}$ on $T$ has kernel exactly
  $T_{\mathrm{spl}}$ and rank $n(n-1)/2$. Each of the seven single classes
  $\gamma$, $\ell$, $\gamma+\ell$, $\gamma+2\ell$, $\beta$, $\widehat\beta$,
  $\beta+\ell$ already has $\{v\in T:v\lrcorner\theta=0\}=T_{\mathrm{spl}}$,
  and each of $150$ random rational classes off $\QQ\eta$ is contracted
  nontrivially by some vector of $T$, while $\eta$ and $3\eta$ are contracted
  to zero by all of $T$; this is the mechanism of \Cref{thm:linebundle}. For
  the six line bundles of \Cref{prop:explicitobject} at $n=3$ the computation
  finds kernel $T_{\mathrm{spl}}$ of dimension $6$ and rank $3$, and checks
  that $\sigma_{E}$ kills the image. It computes the annihilator of
  $\omega_{1}$ in $H^{0,2}$, of dimension $9$, checks that $\sigma_{E}$ kills
  $z\cdot\id_{E}$ for $z$ in it, and finds that the sum of that
  nine-dimensional space and the three-dimensional image has dimension $12$.
  Finally it recomputes the kernel of $\sigma_{E}$ as $15$, and checks that
  the components of every obstruction vector on $L_{i}$ and on
  $L_{i}^{-1}[3]$ are negatives of each other.
\item \emph{The evaluation map.} The computation rewrites the model of item
  (XIII) in its Hodge basis and checks \Cref{prop:evaluation} for the
  explicit object at $n=3$. It builds a basis of $HT^{2}$, of dimension
  $15+36+15=66$, implements the three actions of \eqref{eq:HT2} as
  multiplication, a derivation and a double interior product, and verifies
  the identity $\sigma_{E}(\operatorname{ev}_{E}(x))=x\lrcorner\ch(E)$ on
  every basis vector and on random elements. It confirms that $\ch(E)$ is
  concentrated in degree six and proportional to $\omega_{1}$, and that
  $\operatorname{ev}_{E}$ has rank $45$, against $\dim\Ext^{2}(E,E)=90$. The
  three kernels of contraction with $\ch(E)$ have dimensions $9$, $18$ and
  $9$; their images under $\operatorname{ev}_{E}$ have dimensions $9$, $6$
  and $9$, the first and third coincide, the first and second meet in zero,
  and the sum of all three has dimension $15$, equal to the dimension of
  $\ker\sigma_{E}$, which is recomputed. The last comparison is repeated for
  every one of the eighteen configurations of item (XIV), with
  $\dim\ker\sigma_{E}=15$ and
  $\dim(\ker\sigma_{E}\cap\image\operatorname{ev}_{E})=15$ in every case.
\item \emph{Markman's candidate in dimension eight.} The computation works in
  $\QQ[\Theta]/(\Theta^{5})$ with $\int\Theta^{4}=24$ and checks
  \Cref{prop:markman8} for every odd squarefree $d\le101$: the Euler
  characteristics $14$, $-36$ and $24$ from the Koszul complexes; the count
  $12N(N-5)=(d+9)(3d-3)$ of isolated points and its comparison with the
  number $(d+9)(2d-1)$ of those that are normalised;
  $\chi(\cO_{\widetilde Z})=(d+9)(27-d)$ from the
  Mayer--Vietoris complex; the equality $\ch(\cF)=u_{t}+3v_{t}$ degree by
  degree for the twist $3\Theta$, and the failure by $-4\Theta^{2}$ in
  degree two for the twist $\Theta$; $\Nm(3+\sqrt{-d})=2N$ and the
  discriminant $(d+9)\Theta^{2}$; and the ranks $8d(d-9)$ and $8d(d+9)$ of
  the two transforms, with the threefold control $-8q$ that reproduces
  Markman's rank at $n=3$.
\item \emph{The classes preserving a secant Chern character.} The
  computation works in the Hodge algebra of an abelian $n$-fold with
  $t=\sum_{a}p_{a}\wedge q_{a}$, modulo a large prime,
  and checks \Cref{thm:factor}(i), (ii) and \Cref{lem:bfield}(i) for
  $(n,d,a,b)$ equal to $(3,3,1,1)$, $(3,3,1,3)$, $(4,3,1,3)$, $(4,5,1,3)$,
  $(4,7,2,1)$ and $(5,3,1,3)$, the third being Markman's candidate. It builds
  $\ch(F)=au_{t}+bv_{t}$, implements the actions of $HT^{1}$ and $HT^{2}$ as
  multiplication, interior product, derivation and double interior product,
  and finds: that the kernel of $y\mapsto y\lrcorner\ch(F)$ on $HT^{1}$, a
  space of dimension $2n$, is zero in every case; that the kernel on $HT^{2}$ has
  dimension exactly $n^{2}$, contains the $n(n+1)/2$ polarised deformations
  and the $n(n-1)/2$ classes $(\tfrac d2\,\pi\lrcorner t^{2},0,\pi)$, and is
  spanned by them; and, as controls, that the bare classes $(0,0,\pi)$ and
  the classes $(-\tfrac d2\,\pi\lrcorner t^{2},0,\pi)$ with the compensation
  of the wrong sign all act nontrivially on $\ch(F)$. It then verifies the
  identity $x\lrcorner(we^{B})=((e^{B}x)\lrcorner w)e^{B}$ of
  \eqref{eq:bfield} on random even classes $w$, random $x\in HT^{2}$ and
  $B=kt$ for $n=3,4$, and the count $n(n+1)/2+n(n-1)/2=n^{2}$ for $n<40$.
\item \emph{The local products at a double point.} The computation recomputes
  \Cref{lem:localproducts} over $\QQ[x_{1},\dots,x_{4}]$. For each of the
  two local models, the ideal $I=(x_{1},x_{2})\cap(x_{3},x_{4})$ and the
  complex $K=[R\to R/(x_{1},x_{2})\oplus R/(x_{3},x_{4})]$, resolved by the
  cone of the lifted map, it builds the jet chain maps $-\partial_{a}(d)$ on
  the free resolution, checks that they anticommute with the differential
  and are not null homotopic, and for every pair $(a,b)$ decides whether the
  composite is null homotopic. It finds $\Ext^{2}=k^{4}$ for $I$, with the
  four mixed products nonzero and spanning it and the same-plane products
  zero; $\Ext^{2}=k^{2}$ for $K$, with the two same-plane products nonzero
  and spanning it and the mixed products zero; anticommutation
  $a_{a}a_{b}+a_{b}a_{a}=0$ in both models; and, for $I$, the same vanishing
  pattern from the jet modules and their Yoneda products.

\item \emph{A rank one secant object with smooth support.} The computation
  expands $\ch(\cO_{S})=1-\ch(\cF)e^{-b\Theta}$ directly in
  $\QQ[\Theta]/(\Theta^{5})$, for every $b\le6$ and every squarefree $d<60$,
  and reads off the class $N\Theta^{2}$ with
  $N=\Nm_{K/\QQ}(b+\sqrt{-d})/2$, the pushforward
  $\iota_{*}K_{S}=(4bN/3)\Theta^{3}$ and the Euler characteristic
  $\chi(\cO_{S})=4N(2b^{2}-N)$, checking each against the closed form rather
  than quoting it. It then verifies that Noether's formula and the
  self-intersection formula together give $K_{S}^{2}=12N(4b^{2}-N)$ and
  $e(S)=12N(4b^{2}-3N)$, that the Hodge index bound $9N\ge4b^{2}$ is implied
  by $d>0$ and so never binds, that $e(S)\ge0$ is $d\le\tfrac53b^{2}$ and that
  $3e(S)-K_{S}^{2}=24(b^{4}-d^{2})\ge0$ is $d\le b^{2}$, and that equality in
  the Hodge index would need $9N=4b^{2}$, which no admissible pair reaches.
  At $b=3$ it lists the six surviving discriminants, the four among them with
  $N$ an integer, and the entries of \Cref{tab:smoothsupport}. This is
  \Cref{thm:smoothinvariants} and \Cref{cor:fourdiscriminants}.

\item \emph{The germs that carry the local obstruction.} For ten
  codimension two germs in $\CC^{4}$ the computation determines the number of
  minimal generators, the projective dimension of $R/I$, the ranks of a
  minimal free resolution, and $\cE xt^{1}$, $\cE xt^{2}$ and $\cE xt^{3}$ of
  the ideal against itself; for the germs with $\cE xt^{2}\ne0$ it also
  decides, by null homotopy of the composite of the two jet chain maps, which
  products of translation classes are nonzero. The seven Cohen--Macaulay germs, namely the smooth one, two
  and three planes meeting along a line, the double plane, the quadric cone,
  the fat plane and the cone over the twisted cubic, all have
  $\cE xt^{2}=0$, the last two despite needing three generators in codimension
  two. The three that are not Cohen--Macaulay, namely two planes meeting at a
  point, three planes meeting pairwise at a point and a plane with an embedded
  point, have $\cE xt^{2}\ne0$ and nonzero products. This is
  \Cref{thm:cmvanishing} and \Cref{rem:cmsharp}.

\item \emph{Cohen--Macaulay supports with a split resolution.} The
  computation reads the coefficients $(1,b,-d,-db,d^{2})$ off $u_{\Theta}+bv_{\Theta}$ rather than quoting them,
  confirms the three identities $c_{k+2}+dc_{k}=0$ that make the weighted
  measures of \Cref{thm:noci} share mass, mean and variance, and checks that
  the quadratic of that theorem has discriminant $-\tfrac29b^{2}-2d<0$ for
  every $b\le11$ and every squarefree $d<200$, so that no complete
  intersection can occur. It then enumerates \eqref{eq:burchsystem} exactly,
  by meet in the middle on the four moments, over every multiset of twists in
  bounded boxes, imposing the two conditions of
  \Cref{lem:burchmoments}. Nothing survives at $r=2$ or $r=3$, for $b=3$ and
  every squarefree $d<24$ and also for every $b\le6$ over seven fields; at
  $r=4,5,6$ only $d=15$ and $d=23$ survive, and the first candidate at each
  rank is recorded. The moment system alone does not suffice: at $r=4$,
  $d=23$ it has the solution $p=(-8,-5,0,1,2)$, $q=(-7,-6,-4,4)$, which only
  the injectivity of $\varphi$ removes. This search is not used in the
  paper, which proves \Cref{prop:lowrankburch} and \Cref{thm:nosplit} by
  hand.

\item \emph{The closure graph.} The computation carries the rule set of
  \Cref{setup:rules} as data: forty-two statements, each marked proved here,
  quoted, open or derived, and twenty-two rules, each carrying the label of
  the theorem that proves it. It checks that every statement named in a rule
  is declared and every declared statement is used, that the rule set is
  acyclic, and that each of the seventeen theorem labels it names, for its
  rules and for the consequences of $\mathbf{HC}$ it records, is a label of
  this paper, against the list of all labels generated from the
  sources. It then computes the forward closure by iterating the one step map
  to a fixed point, and verifies the following. No statement marked proved
  rests at any depth on an open one, and $\mathbf{HC}$ is not in the closure
  of the proved and quoted statements. By exhaustive search over all
  $2^{16}$ subsets of the sixteen open statements, the minimal sufficient
  sets are exactly the five of \Cref{thm:closure}(ii), $\{\mathrm{F3}\}$,
  $\{\mathrm{L},\mathrm{M}\}$, $\{\mathrm{L},\mathrm{F3}'\}$,
  $\{\mathrm{V},\mathrm{F3}'\}$ and
  $\{\mathrm{P2}_{\mathrm{split}},\mathrm{F2},\mathrm{F3}'\}$, with
  $\mathrm{P2}_{\mathrm{split}}$ the propagation statement for the split
  families of the CM fields of degree at least four. The statement (F3) is
  the one open statement that suffices alone, and \textup{(F3$'$)} alone
  gives neither the conjecture for abelian varieties nor $\mathbf{HC}$. The
  sets
  $\{\mathrm{P2}_{\mathrm{split}},\mathrm{F2},\mathrm{F3}\}$,
  $\{\mathrm{L},\mathrm{F3}\}$ and $\{\mathrm{V},\mathrm{F3}\}$ suffice and
  are not minimal. The route through this paper is the only minimal set with
  three elements and the only one whose derivation uses a statement proved
  here. Each element of each minimal set is necessary, every minimal set
  contains (F3), \textup{(F3$'$)} or \textup{(M)}, every element of every
  minimal set other than $\mathrm{P2}_{\mathrm{split}}$ is recorded as a
  consequence of $\mathbf{HC}$, and $\{\mathrm{F3}\}$ and
  $\{\mathrm{L},\mathrm{M}\}$ are each equivalent to $\mathbf{HC}$.
  The statement \textup{(V)} alone gives the conjecture for abelian varieties
  and not $\mathbf{HC}$, and \textup{(L)} gives \textup{(V)}. The six open
  demands of the secant route lie outside every minimal set, and granting the
  two of
  them that matter reaches $\mathbf{W}(K,n,\delta_{0})$ and, through the
  descent rule of \Cref{prop:descent}, $\mathbf{W}(K,n,\delta)$ for every
  $\delta$, but not the Weil classes of the fields of degree at least four.
  Within the rule set, the algebraicity of the Weil classes of trivial
  discriminant yields the whole imaginary quadratic case, and so does that
  of the Weil classes of the split families of the fields of degree at least
  four; the propagation statement for the split imaginary quadratic families
  yields all of that case, and $\mathrm{P2}_{\mathrm{split}}$ yields the
  algebraicity of the Weil classes of every CM field, which is therefore a
  derived statement of the rule set and not an open one.
  No minimal set contains a propagation statement for a family of nontrivial
  discriminant or for an imaginary quadratic field, and the bounded
  criterion is absent from the rule set, being equivalent to its own
  conclusion. The computation then removes the rule of \Cref{prop:f3ishc}
  and the statement \textup{(F3$'$)} with its rule and repeats the search,
  which returns the four sets obtained without them,
  $\{\mathrm{L},\mathrm{M}\}$, $\{\mathrm{L},\mathrm{F3}\}$,
  $\{\mathrm{V},\mathrm{F3}\}$ and
  $\{\mathrm{P2}_{\mathrm{split}},\mathrm{F2},\mathrm{F3}\}$; so these two
  additions account for the whole difference between the two answers. It
  also adjoins the rule $\mathrm{F2}\Rightarrow$ (the conjecture for abelian
  varieties) of \Cref{prop:f2isab}, which the rule set does not contain, and
  finds five minimal sets again, $\{\mathrm{F3}\}$,
  $\{\mathrm{L},\mathrm{M}\}$, $\{\mathrm{L},\mathrm{F3}'\}$,
  $\{\mathrm{V},\mathrm{F3}'\}$ and $\{\mathrm{F2},\mathrm{F3}'\}$, none
  of which contains $\mathrm{P2}_{\mathrm{split}}$ (\Cref{rem:f2rule}).
  Finally it checks that the numbering under which the same graph is carried
  into Section 24 of the Lean certificate maps these rules onto the rules
  there. The full derivation of $\mathbf{HC}$ from the proved and quoted
  statements and each minimal sufficient set is written out, one rule per
  line. This is \Cref{thm:closure}, \Cref{prop:f3ishc}, \Cref{prop:f3prime}
  and \Cref{cor:fourremain}, and the count of \Cref{rem:f2rule}.
\item \emph{No split resolution satisfies the weakened criterion.} The
  computation records the data behind \Cref{thm:nosplit}: that the scalar
  $(e^{2}+d)/2$ by which the compensated Poisson class acts on
  $\cO(e\Theta)$ is nonzero for every twist and every $d>0$ in a range; the
  vanishing pattern $H^{2}(\cO(m\Theta))=0$ for $m\ne0$ on an $n$-fold with
  $3\le n\le6$, from the index theorem; that the three split resolutions of
  item (XXXII) have no twist in common between $E_{0}$ and $E_{1}$ and
  satisfy the moment identities \eqref{eq:burchsystem}, so that the theorem
  excludes them as they stand; and, as a control, that no Koszul resolution
  of a complete intersection, twisted by any $b\in[-6,12]$, is secant for
  $d<30$, which is \Cref{thm:noci} in that range.
\item \emph{Two statements that are not reductions, and the consequences of
  injectivity.} The computation has three parts. First, for an abelian
  fourfold $X$ of Mumford type it decomposes every $\bigwedge^{k}V$,
  $V=V_{1}\otimes V_{2}\otimes V_{3}$, under
  $\mathrm{SL}_{2}^{3}$ by peeling highest weights, finds one invariant in
  $\bigwedge^{2}V$ and one in $\bigwedge^{4}V$, counts eight invariants in
  $\bigwedge^{4}(V\oplus V)$ distributed $1,1,4,1,1$ over the K\"unneth
  summands and three in $\bigwedge^{2}(V\oplus V)$, multiplies the three
  divisor classes $\psi_{11},\psi_{12},\psi_{22}$ in the exterior algebra of
  $V\oplus V$ and finds their six products independent, and counts the
  $\Sp_{8}$-invariants of $\bigwedge^{4}(V\oplus V)$ by the Weyl character
  formula over the $384$ signed permutations, finding six. This is
  \Cref{prop:mumfordproduct}. Second, for the quintic threefold it computes
  the Hodge numbers $(1,101,101,1)$ from the Jacobian ring of the Fermat
  quintic, checks in a four-dimensional model that the map from the $(0,3)$
  part to the $(3,0)$ part is symplectic of adjoint weight $3$, and lists the
  adjoint weights, $\{-1,0,1\}$ in weight one against $\{-3,\dots,3\}$ here.
  This is \Cref{prop:noabeliantype}. Third, for $n=2$ and $n=3$ it builds
  $H^{*}(A,\CC)$ with the eigenspace decomposition, the Weil class
  $\alpha_{+}+\alpha_{-}$, and the action of all of $HT^{2}$ on it, and finds
  the annihilator, of dimension $4n^{2}$, spanned by the mixed block, the
  $K$-linear deformations and the mixed bivectors, the remaining classes
  acting injectively. It computes the ideal generated by the mixed block in
  the exterior algebra on $2n$ generators, of codimension $2^{n+1}-1$ with
  quotient of dimension $2\binom{n}{k}$ in degree $k\ge1$, and checks over
  $\mathbb{F}_{3}$ in dimensions two and three that two subspaces spanning
  the space with all mixed wedge products zero cannot both be nonzero. This
  is \Cref{thm:p2forces}.
\item \emph{Weil classes of CM fields of degrees four and six.} The
  computation builds the cyclotomic fields $\QQ(\zeta_{5})$,
  $\QQ(\zeta_{7})$, $\QQ(\zeta_{8})$ and $\QQ(\zeta_{9})$ exactly, as
  $\QQ[x]/(p)$ with elements represented by rational coordinate vectors, and
  carries out the construction of \Cref{thm:cmbasepoint} for six pairs
  $(F,n)$, with $m=2$ and $m=3$ and $n=1,2,3$. For each it checks that
  the CM type and its conjugate partition the embeddings, that
  $\zeta-\zeta^{-1}$ is totally imaginary, that the trace-dual bases satisfy
  $\sum_{k}\sigma(\omega_{k})\sigma'(\omega^{*}_{k})=\delta_{\sigma\sigma'}$
  over every pair of embeddings, that $V^{1,0}$ is half-dimensional and
  isotropic for $E=\Tr(\xi H)$ so that the point lies in the period domain,
  that every eigenspace meets $V^{1,0}$ in dimension $n$ so that the Weil
  classes are of type $(n,n)$, that
  $\delta_{i}(f)=\sum_{\sigma}\sigma(f)\,u_{i,\sigma}\wedge u_{i',\sigma}$
  and is therefore of type $(1,1)$, that every $\alpha_{\sigma}$ is nonzero,
  and that the balanced $n$-fold product of the $\delta_{i}$ is
  $w(f)=\sum_{\sigma}\sigma(f)\alpha_{\sigma}$, rational, spanning $W_{F}$.
  For $F=\QQ(\zeta_{5})$ at $n=2$ it also computes, in the exterior algebra of
  the sixteen-dimensional space, the space of rational $(1,1)$ classes, of
  dimension $32$, and records that the balanced sum is necessary: the plain
  product $\delta_{1}(1)\delta_{2}(1)$ is not a Weil class. Finally, for
  $F=\QQ(\zeta_{8})\supset K=\QQ(\sqrt{-1})$ it checks that each $K$-Weil class
  is the product of the two $F$-Weil classes above it and that $W_{K}$ lies in
  the span of the products of pairs from $W_{F}$, which is
  \Cref{prop:composite}.
\item \emph{Transport along the orbit.} The computation works on the
  lattice $\Lambda=\ZZ^{2G}$ with the standard alternating form. For a sample
  of rational symplectic elements drawn from the unipotent radical of the
  stabiliser of an isotropic line, with denominators up to $29$, it verifies
  the three scaling identities of \Cref{lem:transport}, namely
  $\det\varphi=c^{2G}$, $\varphi^{*}E=c^{2}E$ and
  $\bigwedge^{k}\varphi=c^{k}\bigwedge^{k}g$, the last on the exterior
  algebra for $G=2,3,4$, and records that the denominators are unbounded. It
  then takes two sublattices of rank $2n=4$ in an abelian fourfold, one
  preserved by the sample elements and one not, computes the multiplicity
  $m=[\varphi^{-1}(\Lambda)\cap W:\Lambda_{W}]$ by Smith normal form and the
  degrees as Pfaffians, and checks
  $\deg\varphi(B)=c^{2n}\deg(B)/m$ in every case. The two cases separate: for
  the preserved subtorus $m=c^{2n}$ exactly and the image degree is constant,
  while for the other $m=c^{2n-1}$ and the image degree grows linearly in $c$.
  Finally it tabulates the bound $c^{2n}\deg(Z)/|G(Z)|$ of
  \Cref{thm:p1forces}(ii) for stabilisers of order $1$, $2$ and $4$, which is
  strictly increasing and unbounded. This is \Cref{lem:transport},
  \Cref{lem:multiplicity} and \Cref{thm:p1forces}.
\item \emph{The annihilator of a Weil class in Hochschild cohomology.} The
  computation builds $HH^{*}(A)=\bigwedge^{*}\bigl(H^{0,1}\oplus
  H^{0}(T_{A})\bigr)$, of
  dimension $2^{4n}$, as an exterior algebra on $4n$ operators, with a basis
  indexed by subsets of them, acting on $H^{*}(A,\CC)$ by wedging and
  contracting, and computes the annihilator of
  $\omega=a\alpha_{+}+b\alpha_{-}$ in every degree by exact rational
  elimination. In degree one, $P=V_{+}^{0,1}\oplus T_{-}$ is exactly the
  annihilator of $\alpha_{+}$ and $Q=V_{-}^{0,1}\oplus T_{+}$ that of
  $\alpha_{-}$, and $\Ann_{HH^{1}}(\omega)=0$. In degree two,
  $\Ann_{HH^{2}}(\omega)=P\wedge Q$, of dimension $4n^{2}$. In every degree
  $k\ge1$ the annihilator has dimension
  $\dim\Ann_{HH^{k}}(\omega)=\binom{4n}{k}-2\binom{2n}{k}+\delta_{k,2n}$, and
  in degree $0$ it is $0$. The computation also confirms that
  the ideal generated by $P\wedge Q$ consists of the monomials meeting both
  halves and has codimension $2^{2n+1}-1$, that contraction with
  $\alpha_{-}$ is injective on $\bigwedge^{*}P$ and contraction with
  $\alpha_{+}$ on $\bigwedge^{*}Q$, the two images meeting in exactly one
  line in degree $2n$, and that $\dim HH^{2}-\dim\Ann_{HH^{2}}=2n(2n-1)$. It
  is run in every degree for $n=1$, where degree two is the exceptional
  degree $2n$ and the annihilator is one larger, and for $n=2$ with two
  choices of $(a,b)$, and in degrees up to three for $n=3$. This is
  \Cref{thm:hhsplitting} and \Cref{cor:hhfactor}.

\item \emph{The numerical criterion and the support of an object meeting
  it.} The computation checks the finite linear algebra of
  \Cref{thm:p2numerical}, \Cref{lem:finiteproduct} and \Cref{thm:p2support}
  in five parts. (A) In the exterior algebra of
  $H^{*}(A,\CC)$ with $HH^{1}$ acting by wedging and contracting, the map
  $\xi\mapsto\xi\lrcorner\,\omega$ on $HH^{2}$ has rank exactly $2n(2n-1)$
  for $n=2$ and $n=3$, that is $12$ and $30$, and it is injective on
  $\bigwedge^{2}P\oplus\bigwedge^{2}Q$; and
  $\binom{4n}{2}-4n^{2}=2\binom{2n}{2}=2n(2n-1)$ for every $n\le60$. (B) On a
  torus of dimension $c$, for $2\le c\le6$, the class of a point contracted
  with $a\wedge b$ is nonzero for all coordinate pairs and for twenty-five
  random rational pairs of independent vectors, and zero for twenty-five
  dependent pairs. (C) For $n=2,3,4$ and random subspaces $W$ of
  $T_{+}\oplus T_{-}$ of codimension at least two, together with the boundary
  cases where $W$ contains $T_{+}$ or $T_{-}$, the images of $T_{+}$ and
  $T_{-}$ in the quotient have all mixed wedges zero exactly when $W$
  contains $T_{+}$ or $T_{-}$; this is the dichotomy used in the proof of
  \Cref{thm:p2support}(i). (D) In an explicit rational model of $\RR^{4n}$
  with complex structure $J$ and an action $M$ of $\QQ(i)$ of signature
  $(n,n)$, scrambled by a random rational change of basis, the forms
  vanishing on $\ker(J-M)$, the real form of $T_{+}$, are exactly
  $V_{-}^{1,0}\oplus V_{+}^{0,1}$, computed over the Gaussian rationals for
  $n=1,2$. (E) For $2\le n\le6$ neither $\alpha_{+}$ nor $\alpha_{-}$ is a
  monomial of $\bigwedge^{*}(V_{-}^{1,0}\oplus V_{+}^{0,1})$ or of
  $\bigwedge^{*}(V_{+}^{1,0}\oplus V_{-}^{0,1})$, so the Weil line meets the
  sum of the two subrings in zero.

\item \emph{Two jet classes on a complex of finite length.} For a bounded
  complex $M$ of free modules over $\QQ[x_{1},x_{2}]$ or
  $\QQ[x_{1},x_{2},x_{3}]$ with finite length cohomology, the computation
  builds the jet chain maps $-\sum_{i}a_{i}\partial_{i}(d)$ and decides
  whether the composite of two of them is null homotopic. It first confirms
  that the cohomology of each input has finite length, and then tests every
  pair of coordinate directions and a random pair of independent rational
  directions on eighteen modules: in two
  variables the residue field, $R/\mathfrak{m}^{2}$, $R/\mathfrak{m}^{3}$,
  four monomial ideals, the complete intersection
  $(x_{1}^{2}+x_{2}^{2},x_{1}x_{2})$, the ideal
  $(x_{1}^{3}-x_{2}^{3},x_{1}^{2}x_{2},x_{1}x_{2}^{2})$, the non-cyclic
  modules $\mathfrak{m}/\mathfrak{m}^{3}$ and the cokernel of
  $\left(\begin{smallmatrix}x_{1}&x_{2}&0\\0&x_{1}&x_{2}\end{smallmatrix}\right)$,
  and $k\oplus R/\mathfrak{m}^{2}$; in three variables the residue field,
  $R/\mathfrak{m}^{2}$, two monomial complete intersections, a monomial
  ideal of colength five and the Gorenstein algebra
  $R/(y_{1}^{2}-y_{2}^{2},y_{2}^{2}-y_{3}^{2},y_{1}y_{2},y_{2}y_{3},y_{1}y_{3})$;
  and three complexes with several cohomology modules. In every case the
  Euler characteristic is nonzero and every product of two independent jet
  classes is nonzero, which is \Cref{lem:finiteproduct}. It then builds the
  cone $C$ of $J(\partial_{2})J(\partial_{1})\colon K\to K[2]$ on the Koszul
  complex $K$ of the residue field in two variables, checks that $C$ is a
  complex with cohomology of lengths $1$ and $1$ and Euler characteristic
  zero, and that $J(\partial_{2})J(\partial_{1})$ is null homotopic on $C$,
  while on $k\oplus k[1]$, also of Euler characteristic zero, it is not. So
  the hypothesis $\chi(M)\ne0$ of the lemma cannot be dropped.

\item \emph{The two classes on a Mumford fourfold, and the K3 surface that
  would carry them.} The computation works on the split model
  $V=V_{1}\otimes V_{2}\otimes V_{3}$ of
  \Cref{setup:mumford}, with $\mathfrak{sl}_{2}^{3}$ acting through its nine
  generators and $\psi=\epsilon\otimes\epsilon\otimes\epsilon$, in five parts.
  (A) The commutant of $\mathfrak{sl}_{2}^{3}$ in $\End\bigwedge^{2}V$,
  computed as the null space of a linear system in $784$ unknowns, has
  dimension four, and the four projections built from the three Casimir
  operators are orthogonal idempotents of ranks $1$, $9$, $9$, $9$ that sum to
  the identity, commute with the action and span the commutant. (B) The Hodge
  numbers are $(6,16,6)$ on $H^{2}$, $(0,9,0)$ on $U_{12}$, $(3,3,3)$ on
  $U_{13}$ and $U_{23}$, and $(1,7,1)$ on $T(-1)$. (C) The product map $\mu$ of
  \Cref{lem:mumfordmu} carries $T_{i}\otimes T_{j}$ onto $U_{ij}$ with rank nine
  for $i\ne j$ and $T_{i}\otimes T_{i}$ onto the line of $\theta$, and its
  image is all of $H^{2}$. (D) The Lefschetz pairing
  $\int\alpha\beta\theta^{2}$ is nondegenerate on $H^{2}$; $\mu\mu^{\dagger}$
  acts on the three $U_{ij}$ by one and the same nonzero rational number
  $\nu$, which is $-1$ in the normalisation of the model; for three choices of
  the conjugates of $e$ and $\lambda$, $\mu_{e}\mu_{e}^{\dagger}$ acts on
  $U_{ij}$ by $\nu\sigma_{i}(g)\sigma_{j}(g)$ with $g=e^{2}/\lambda$ and
  preserves the line of $\theta$; the algebra it generates on $U$ is spanned by
  the three projections when the conjugates of $g$ are distinct and is the
  scalars when they coincide; and $\kappa^{\dagger}\kappa$, for the trace form
  on $\End(V)$, is multiplication by $4e^{2}/\lambda$. (E) In the Clifford
  algebra of $T$ with the form $\sum_{i}l_{i}\operatorname{tr}(x_{i}y_{i})$,
  $(l_{1},l_{2},l_{3})=(1,2,5)$, with a basis indexed by subsets of an
  orthogonal basis of $T$, the spin lifts of the three copies of
  $\mathfrak{sl}_{2}$ act on $C^{+}(T)$, of dimension $256$; the vectors of highest weight $(1,1,1)$ span $32$
  dimensions and lowering them gives a basis, so $C^{+}(T)\cong V^{\oplus32}$.
  In that basis the map $\Psi(v)x=vxv_{0}$ of \Cref{prop:ksembed} is
  $v\otimes N_{i}$ for $v\in T_{i}$, the three matrices $N_{i}$ anticommute
  pairwise, square to the scalars $4$, $8$, $20$, proportional to the $l_{i}$,
  and are linearly independent. This is \Cref{prop:mumfordrm},
  \Cref{lem:mumfordmu} and the shape of the Kuga--Satake map used in the proof
  of \Cref{thm:mumfordks}.

\item \emph{The Lefschetz operator of an abelian scheme over a curve.} The
  computation works in exact exterior algebra over $\QQ$ and assumes no sign
  convention. (A) On ten abelian varieties with $g\le4$ and polarisation
  types from $(1)$ to $(1,2,4)$ and
  $(1,1,2,2)$, the contraction by the dual bivector satisfies
  $[L,\Lambda]=H$, and the Pontryagin product, computed from
  $x\star y=\mu_{*}(p_{1}^{*}x\cup p_{2}^{*}y)$ with $\mu_{*}$ defined against
  $\mu^{*}$ by the projection formula, gives $D^{-1}(x\star\gamma)=\Lambda x$
  in every degree, with $\ell\star\gamma=gD$; without the factor $D^{-1}$ the
  identity fails at type $(1,2)$. (B) On four product families $E\times A$
  over an elliptic curve, where every class is invariant, the operator
  $D^{-1}T+\Lambda_{C}$ of the proof of \Cref{thm:lefschetzfamily}, with
  $\Lambda_{C}$ built from bases of invariant classes and the inverse of their
  pairing matrix, satisfies $[L_{W},\Lambda]=H$, and the cross relations
  $[L_{\mathrm{rel}},\Lambda_{C}]=0$ and $[L_{C},T]=0$ hold. (C) The invariants
  of $\mathfrak{sl}_{2}^{3}$ in $\bigwedge^{q}(V_{1}\otimes V_{2}\otimes V_{3})$,
  computed as the joint kernel of the six generators, have dimensions
  $1,0,1,0,1,0,1,0,1$. This is \Cref{prop:pontryaginlambda}, the algebra of
  \Cref{thm:lefschetzfamily} and the count in \Cref{cor:mumfordlefschetz}.

\item \emph{Reduction of the two targets.} The computation works in the
  split model $H^{*}(X\times X)=\bigwedge^{*}(V\oplus V)$, on sixteen
  generators, in exact rational arithmetic where a space is to be
  determined, and with ranks of integer vectors modulo the prime $1000003$
  where a span is to be shown full, a rank modulo a prime being a lower bound
  for the rank over $\QQ$. (A) The invariants of $\mathfrak{sl}_{2}^{3}$ in $\bigwedge^{k}(V\oplus V)$ have
  dimensions $1,0,3,0,8,0,16,0,28,0,16,0,8,0,3,0,1$; in degree four the
  products of divisor classes give six, so two are exceptional; and the ring
  of invariants is generated in degrees two and four. (B) At a CM point the
  Hodge classes of $X_{c}$ number $1,0,4,0,8,0,4,0,1$ and those of
  $X_{c}\times X_{c}$ number $1,0,16,0,132,0,432,0,648,\dots$; the divisor
  classes of $X_{c}\times X_{c}$ span only $100$ of the $132$ in degree four,
  while the ring generated by them and by the pull-backs of the Hodge classes of
  $X_{c}$ along $(\varphi,\psi)\colon X_{c}\times X_{c}\to X_{c}$ spans every
  Hodge class in every degree. (C) On an abelian $2n$-fold of Weil type,
  $n=1,2,3$, the classes of type $(p,p)$ annihilated by $Q\wedge HH^{1}$, by
  $P\wedge HH^{1}$ and by $HH^{1}\wedge HH^{1}$ are $\CC\alpha_{-}$,
  $\CC\alpha_{+}$ and $0$. This is \Cref{thm:mumfordcm}, the ring structure
  used in \Cref{cor:mumfordequiv}, and the annihilators in the proof of
  \Cref{thm:nosum}.

\item \emph{Rigidity of the exceptional classes of a Mumford square.} The
  computation works in the model of \Cref{setup:omegaa} at a point of the
  Mumford curve, with $V^{1,0}$ the vectors of third weight $+1$, in exact
  rational arithmetic with the coefficients $a_{0},\dots,a_{3}$ as symbols.
  (A) The Casimir operators split ${\textstyle\bigwedge^{2}}V$ as $1+9+9+9$,
  and the four tensors $\pi$ are invariant. (B) $\mathfrak{sp}^{-1,1}$ has dimension $36$ and splits into nine weight
  blocks of dimensions $3,4,3,4,8,4,3,4,3$; the generic kernel of contraction
  into $\omega_{a}$ is one vector, independent of $a$; every Gram determinant
  is a positive constant times a product of $P_{1},P_{2},Q_{1},\dots,Q_{4},R$;
  each is a sum of squares whose zeros force two of $a_{1},a_{2},a_{3}$ to
  coincide; the tangent dimensions at $(3,5,-7,11)$, $(3,5,5,5)$ and
  $(1,1,1,1)$ are $1$, $10$ and $20$. (C) $HT^{2}=28+64+28$ splits into $46$
  blocks by type, weight and exchange parity, $\omega_{a}$ being exchange
  invariant; every Gram determinant is a positive constant times a product of
  the forms of (B), $S_{1},\dots,S_{4}$, $Z_{1},\dots,Z_{6}$, $a_{2}$, $a_{3}$
  and four further linear forms; the only block with a kernel at a generic
  point is the even block of weight zero in $H^{1}(T)$, and its kernel is the
  lowering derivation of the third factor on both copies; $L_{0}$ divides only
  the nine odd bivector blocks. (D) Certificates: each $S$ is a sum of squares
  whose zeros force a coincidence or a zero among $a_{1},a_{2},a_{3}$, each
  $Z$ is a positive combination of four independent squares, and the linear
  factors other than $L_{0}$ have unequal coefficients on $a_{1},a_{2},a_{3}$.
  (E) As a cross-check, the annihilator in $HT^{2}$ has dimension one at forty
  random admissible points, $36$ at $(1,1,1,1)$, $7$ at $a_{2}=0$, $4$ on the
  hyperplane $a_{0}-3a_{1}-3a_{2}+a_{3}=0$, and $10$ at three points of
  $L_{0}=0$. This is \Cref{thm:mumfordrigid}, \Cref{rem:rigidcompare} and
  the annihilator used in \Cref{thm:mumfordnumerical}.

\item \emph{Objects with an exceptional Chern character.} The computation
  works in the same model. (A) At $a=(3,5,-7,11)$ the ranks of contraction
  into $\omega_{a}$ on $HT^{k}$, computed exactly over $\QQ$ block by block, are
  $1,16,119,328,560,328,119,16,1$ for $k=0,\dots,8$ and $0$ above. (B) The
  same ranks hold modulo the prime $1000133$, at which the minimal polynomial
  of $\zeta_{7}+\zeta_{7}^{-1}$ splits, at $a_{0}=2$ and the three conjugates
  in each of the six orders; at a point of $L_{0}=0$ they are
  $1,16,110,328,548,328,110,16,1$. (C) For a random class with components of
  types $(0,0)$ to $(3,3)$ the ranks are palindromic and vanish above degree
  eight. (D) The alternating sum of the ranks is $112$ and
  $2\cdot1+2\cdot119+560=800$. (E) At a CM point there are $16$ divisor
  classes and $132$ Hodge classes in degree four; products of divisor classes
  span $100$, and meet the span of the four $\pi$ exactly where
  $a_{1}=a_{2}=a_{3}$. (F) At a CM point $X_{c}$ has $8$ Hodge classes in
  degree four, of which products of divisor classes span $6$. This is
  \Cref{lem:rankduality}, \Cref{thm:mumfordprofile} and
  \Cref{prop:mumforddivisor}.

\item \emph{The classes generated by line bundles on a Mumford square.} The
  computation works in the same model, with $\mathfrak{sp}(V,\psi)$
  spanned by the $36$ maps $x\mapsto\psi(u,x)v+
  \psi(v,x)u$ and the torus of diagonal matrices in $\Sp(V,\psi)$, of rank
  four. (A) The $36$ maps and the Lie algebra of $G$ lie in
  $\mathfrak{sp}(V,\psi)$. (B) The $\Sp(V,\psi)$-invariants of
  ${\textstyle\bigwedge^{k}}(V\oplus V)$ have dimension at most
  $1,3,6,10,15,10,6,3,1$ for $k=0,2,\dots,16$, as a kernel modulo $1000003$, and products of the three
  invariant divisor classes span spaces of exactly these dimensions. (C) The
  $G$-invariants have dimension $1,3,8,16,28,16,8,3,1$, by weight
  multiplicities and as a kernel, so the Hodge classes that are not
  Lefschetz classes number $0,0,2,6,13,6,2,0,0$, twenty-nine in all. (D) At
  a CM point the monomials of weight zero for the diagonal torus number
  $1,16,100,304,454,304,100,16,1$, the coefficients of $(1+4y+y^{2})^{4}$;
  each is a product of the $16$ divisor classes, and those products span
  exactly them; the Hodge classes number $1,16,132,432,648,432,132,16,1$. (E)
  $\pi_{0}$ and $\pi_{12}+\pi_{13}+\pi_{23}$ are $\Sp(V,\psi)$-invariant and
  $\pi_{12}$ is not; the components of $\pi_{12},\pi_{13},\pi_{23}$ of
  nonzero weight for the diagonal torus have rank two over $\QQ$, with the
  relation $(1,1,1)$ only. (F) The modules generated by $\pi_{12}-\pi_{13}$
  and by $\pi_{13}-\pi_{23}$ have dimension $308$ and highest weight
  $(2,2,0,0)$, and Weyl's formula gives $1,27,42,308$ for the highest weights
  $0,\varpi_{2},\varpi_{4},2\varpi_{2}$, with sum $378$. (G) The module meets
  the span of the four $\pi$ in a plane. This is
  \Cref{thm:lefschetzclosure}.

\item \emph{The Hodge locus of an exceptional class among all complex tori.}
  The computation realises $V_{\RR}$ as the fixed space in $V_{1}\otimes V_{2}\otimes V_{3}$ of the real structure that
  is quaternionic on the first two factors and real on the third, in exact
  Gaussian integer arithmetic. (A) It is a real structure, $V_{\RR}$ has
  dimension $8$, and $\psi$ is real and alternating on it. (B) The complex
  structures $1\otimes1\otimes j_{3}$ of the Mumford family and
  $j\otimes1\otimes1$, $1\otimes j\otimes1$ of the compact factors preserve
  $V_{\RR}$, and $\psi(x,Jy)$ has signature $(0,8)$ for the first and $(4,4)$
  for the other two. The unit quaternion $k\otimes1\otimes1$ with $kj=-jk$
  preserves $V_{\RR}$ and $\psi$ and carries $\psi(x,Jy)$, for
  $J=j\otimes1\otimes1$, to its negative, so every form $\psi\otimes s$ on
  $V_{\RR}\oplus V_{\RR}$ with $s\ne0$ is congruent to its negative and
  indefinite. (C) At $j\otimes1\otimes1$ the tensors $\pi$ have type $(2,2)$ and the annihilator of $\omega_{a}$ in
  $H^{1}(T)$, of dimension $64$, is spanned by the lowering derivation of the
  first factor; for the Mumford family it is spanned by that of the third,
  as in item (XLIV). (D) The exchange of the first and third factors carries
  $\pi_{12}$ to $\pi_{23}$ and fixes $\pi_{0}$ and $\pi_{13}$. This is
  \Cref{prop:notwistor}.

\item \emph{The square as a holomorphic symplectic variety.} The
  computation works in the same split model, with
  $\iota(x)=(x\otimes1)\Psi$ for $x$ in $T=\bigoplus_{i}\mathfrak{sl}(V_{i})$.
  (A) $\iota(x)$ is symmetric and $\iota$ is equivariant, for all pairs of the
  nine generators. (B) The three Casimir operators split
  $\bigwedge^{2}(V\oplus V)$, of dimension $120$, into isotypic pieces of
  highest weights $(0,0,0)$, $(2,2,0)$, $(2,0,2)$, $(0,2,2)$, $(2,2,2)$,
  $(2,0,0)$, $(0,2,0)$, $(0,0,2)$ and dimensions $3,27,27,27,27,3,3,3$, with
  parts of type $(2,0)$ of dimensions $0,0,9,9,9,0,0,1$; the last three pieces
  are spanned by $\iota(T_{1})$, $\iota(T_{2})$, $\iota(T_{3})$. (C) The form
  $\sigma_{0}=\iota(E_{3})$ is of type $(2,0)$ and pairs the vectors of type
  $(1,0)$ of the two copies through an invertible $4\times4$ block. (D)
  $\iota_{2}(C_{1})=3\pi_{0}-\pi_{12}-\pi_{13}+3\pi_{23}$ and cyclically,
  exactly, and hence the formula for $\iota_{2}(\sum s_{i}C_{i})$ with symbolic
  $s_{i}$. (E) $\det(x_{1}\otimes1\otimes1+1\otimes x_{2}\otimes1+
  1\otimes1\otimes x_{3})=\Delta(N_{1},N_{2},N_{3})^{2}$ as a polynomial
  identity in the nine coordinates. (F) $\iota(x)^{8}=8!\det(x)$ times the
  volume form at three random integral points. (G) $\Delta$ is a
  nondegenerate ternary form, of Gram determinant $-4$, and irreducible. This
  is \Cref{prop:hksquare}.

\item \emph{The open routes to the Mumford target.} The computation carries
  the implications into and out of the Mumford target as data, in the manner of item (XXXIII):
  sixteen statements, one of them proved (the Hodge conjecture at the CM
  points), and twenty-two Horn rules, each labelled by the theorem that proves
  it; unlike that rule set, this one contains equivalences. (A) Its eleven
  labels are labels of the paper, and its statements are all used. (B) The
  target is not in the closure of what is proved. (C) The statements equivalent to the
  target over what is proved are exactly the target, algebraicity at
  uncountably many points, $B(W\times_{C}W)$, bounded data at infinitely many
  points, and bounded data modulo $p$. (D) The minimal sets of open statements
  that yield the target, found by exhaustive search, number twelve and each
  has one element. (E) The target yields neither the complex with
  $\dim\Ext^{2}=119$ nor the two Kuga--Satake statements. (F) \textup{(L)}
  and \textup{(V)} each yield the target. (G) The Hodge conjecture and (F2)
  yield the target, and the target yields none of the Hodge conjecture, (F1),
  (F2), (F3), \textup{(L)}, \textup{(V)}, \textup{(M)}. (H) The statement (F3) is
  equivalent to the Hodge conjecture by \Cref{prop:f3ishc}, so it is the
  twelfth one-element route to the target, and the minimal sets of open
  statements that yield the Hodge conjecture in this rule set are the
  conjecture itself, (F3), and \textup{(L)} with \textup{(M)}, none of which
  contains the target. This is \Cref{cor:mumfordone} and
  \Cref{rem:mumfordneeded}.
\item \emph{The shape of an object meeting the numerical criterion.} The
  computation checks the finite ingredients of \Cref{thm:p2endo}. (A) In the
  coordinate model of \Cref{sec:explicit}, for $n=2$ with $d=1,2,3,5,7$ and for $n=3$ with
  $d=1,2,3$, the rational Weil classes $\omega_{1}$, $\omega_{2}$ satisfy
  $\eta\wedge\omega_{1}=\eta\wedge\omega_{2}=0$, and in the orientation in
  which $\eta^{2n}$ is positive the intersection form on the rational Weil
  plane is positive definite at $n=2$ and negative definite at $n=3$, as the
  Hodge--Riemann relations require; its Gram matrix in the basis
  $\omega_{1},\omega_{2}$ is $\mathrm{diag}(8d^{2},8d)$ at $n=2$ and
  $\mathrm{diag}(-32d^{3},-32d^{2})$ at $n=3$. (B) In the model of item
  (XXXVIII), for $n=2$ and $n=3$ and in every degree, the classes of type
  $(p,p)$ killed by all $4n^{2}$ monomials of $P\wedge Q$ are exactly
  $\CC\alpha_{+}\oplus\CC\alpha_{-}$. (C) The generators of
  $\bigwedge^{2n}P$ and $\bigwedge^{2n}Q$ contract $\omega$ to nonzero
  multiples of the generator of $H^{0,2n}$; contraction with $\omega$ is
  injective on $\bigwedge^{2n-2}P\oplus\bigwedge^{2n-2}Q$; and the pairing
  $\bigwedge^{2}P\times\bigwedge^{2n-2}P\to\bigwedge^{2n}P$ is perfect.
  (D) The alternating sum of the lower bounds $1,2\binom{2n}{1},\dots,1$ is
  $-2$ for every $n\le60$; at $n=2$ and $n=3$, exhaustive search over the
  admissible profiles shows $2e_{0}\ge\chi+4$, hence $e_{0}\ge3$, lists the
  smallest profiles, and gives $e_{1}\le e_{0}+9$ at $n=3$; at $n=4$ and
  $n=5$ a profile with $e_{0}=1$ satisfies the same constraints, so the bound
  on $\End$ is special to $n\le3$; and the same identity with $\chi=0$ on the
  Mumford square returns the $800$ odd self-extensions of item (XLV).
\item \emph{The Hodge classes of a very general Weil torus.} The
  computation checks the finite inputs of \Cref{lem:weiltorushodge},
  \Cref{lem:weiltorussubsets} and \Cref{thm:p2false}, in the coordinate model of \Cref{sec:explicit}, for
  $n=2$ with $d=1,3$ and for $n=3$ with $d=2$. The members of the family $S$
  of \Cref{setup:weiltori} used are $T_{X,Y}$ with $X$ and $Y$ spanned by
  explicit $K$-rational combinations of the eigenvectors, three of them for
  $n=2$ and five for $n=3$. (A) The rational Weil classes $\omega_{1}$,
  $\omega_{2}$ are of type $(n,n)$ at every member used. (B) The
  polarisation $\eta$ is of type $(1,1)$ at the balanced member and at none of
  the members used, which therefore lie off the polarised family. (C) The
  rational classes of degree $2k$, $1\le k\le n$, that are of type $(k,k)$
  at all the members used form a space of dimension $0$ for $k<n$ and $2$ for
  $k=n$, the Weil plane. A class of degree $2k$ is of type $(k,k)$ exactly
  when its product with every product of $2n-k+1$ holomorphic $1$-forms
  vanishes; the ranks are computed modulo a prime $q<2^{31}$ in which $-d$ is
  a square, with $\sqrt{-d}$ sent to both square roots so that the conjugate
  conditions are imposed as well; since reduction can only lower a rank, the
  dimensions found are upper bounds, and they are attained by the Weil
  classes. Degrees above $2n$ are not checked; the proof of the lemma covers
  them. (D) For an element $\kappa$ of $V_{+}\otimes V_{-}$ with random
  coefficients, $\alpha_{+}\wedge\kappa^{n}=\alpha_{-}\wedge\kappa^{n}=0$,
  and $\alpha_{+}\wedge(\kappa+\beta)^{n}=\alpha_{+}\wedge\beta^{n}$ becomes
  nonzero once a component $\beta\in\bigwedge^{2}V_{-}$ with
  $\beta^{n}\ne0$ is added, while $\alpha_{-}\wedge(\kappa+\beta)^{n}$ stays
  zero, every term having a factor from $V_{-}$; the same holds with the
  roles of $V_{+}$ and $V_{-}$, and of $\alpha_{+}$ and $\alpha_{-}$,
  exchanged. So the vanishing used for the block diagonal K\"ahler class
  $\kappa_{0}$ is a property of $V_{+}\otimes V_{-}$ and not of the
  coordinates.

\item \emph{The polarised criterion as a number.} The computation works in
  the eigenbasis model of item (XXXVIII), with
  $\theta=i\sum_{j}s_{j}e_{j}\wedge f_{j}$ and
  $\gamma=\sum_{k}c_{k}\theta^{k}+u\alpha_{+}+\bbar u\alpha_{-}$, builds the
  image $\xi\lrcorner\gamma$ of every element of a basis of $HT^{2}$, and
  computes ranks exactly over $\QQ(i)$ by blocks of torus weight. This is
  justified because contraction is equivariant for the torus of
  \Cref{thm:p2primenumber} and elements whose weights differ by other than a
  multiple of $(1,\dots,1)$ have images in different weight spaces; nothing
  else about the structure is used. (A) For $n=3,4$, ten shapes (pure, $c_{0}$
  only, $c_{2n}$ only, $e^{\theta}$, $c_{1}$ only, $c_{0}$ and $c_{1}$, $c_{0}$
  and $c_{2n}$, two exponentials, $c_{n}$ only, and random rational
  coefficients) and two values of $u$: $\dim\Ann_{HT^{2}}(\gamma)=n^{2}(4-\rho)$
  and $r=(4+\rho)n^{2}-2n$, with $\rho$ computed separately as the rank of the
  Hankel matrix. (B) For $n=2,3,4$ and a general shape, the annihilator is
  spanned by the $n^{2}$ vectors $v^{1}_{ab}+v^{2}_{ab}$, lies in $H^{1}(T_{A})$,
  and these vectors kill $\gamma$. (C) $\rho=1$ on $\CC e^{t\theta}$ and
  $\CC\theta^{2n}$, and multiplying $\gamma$ by $e^{t\theta}$ changes neither
  $\rho$ nor $r$. (D) At $n=2$, $\dim\Ann=4(4-\rho)+\dim\ker((HK)^{2}-|u|^{2})$
  on six values of $|u|$ for $\theta^{2}+\omega$ and on four random shapes, and
  $r=24$, $22$ and $23$ at $|u|^{2}=1,4,16$. (E) The first-order Hodge locus
  $\Ann_{H^{1}(T_{A})}(\gamma)$ has dimension $n^{2}$ or, when
  $c_{1}=\dots=c_{2n-1}=0$, $2n^{2}$, for $n=3,4$ and all ten shapes; at $n=2$
  it is $5$ exactly at $c_{1}=c_{3}=0$, $|u|=4|c_{2}|$. (F)
  $\int\ch^{\vee}\ch=(2n)!\,Q(c)+2|u|^{2}$ in the normalisation
  $\int\theta^{2n}=(2n)!$ for $n=2,3$, and in the rational model of item (VI)
  at $n=2$, $\int\omega_{1}^{2}=8d^{2}$ and
  $\int\eta^{4}=24$ for $d=1,2,3,5,7$. (G) At $n=2$ and the exceptional
  ratio $|u|=4|c_{2}|$, the stabiliser of $\omega+c_{2}\theta^{2}$ in
  $\mathfrak{gl}(H^{1})$ has dimension $21$, Hodge pieces $5$, $11$, $5$,
  commutant on $H^{1}$ the scalars, a real form whose trace form has
  signature $(12,9)$, and invariants of dimensions $1,0,0,0,1,0,0,0,1$ in
  $\bigwedge^{k}H^{1}$; at an ordinary ratio the stabiliser has dimension $15$
  and invariants $1,0,1,0,3,0,1,0,1$; at the exceptional ratio
  $\int\gamma\,\kappa_{h}^{2}=\frac16(\sigma_{2}(h)+4\operatorname{Re}\det
  h_{01,23})\int\theta^{4}$ at eight random $h$, and the sign of
  $\int\gamma\,\kappa^{2}$ is that of $c_{2}$ at $150$ random K\"ahler
  classes, while at the ratio $R=3/16$ both signs occur. This is
  \Cref{prop:p2primelocus} and \Cref{rem:p2primeexceptional}. In a longer
  run the closed form of (A) is checked exactly for $n=5,\dots,9$ on five
  shapes, and at $n=10$, where $HT^{2}$ has dimension $780$ and $\gamma$ has
  $2^{20}+2$ terms. This is \Cref{thm:p2primenumber}.

\item \emph{The Hochschild profile of a polarised character.} In the same
  model and by the same torus blocks, the computation determines the rank
  $\rho_{k}$ of contraction into $\gamma$ on every $HT^{k}$, exactly over
  $\QQ(i)$. (A) For $n=2,3$, in every degree $0\le k\le4n$, on eight shapes and two values of $u$, the ranks equal
  $2\binom{2n}{k}+M_{k}\min(k+1,2n+1-k,r')-[k=n]\,d$ for $1\le k\le2n-1$, with
  $\rho_{0}=\rho_{2n}=1$ and $\rho_{k}=0$ above $2n$; the general profiles are
  $(1,8,24,8,1)$ and $(1,12,57,112,57,12,1)$, and the pure case gives
  $2\binom{2n}{k}$. (B) $\rho_{k}=\rho_{2n-k}$ on every shape. (C) For
  $c_{n}\theta^{n}$ alone the middle rank drops by $d$ exactly at
  $|u|=\binom{n}{a}n!|c_{n}|$ and in no other degree. (D) On
  $\bigwedge^{2n}\ZZ^{4n}$, $n\le3$, the form $\int x\wedge y$ pairs each
  monomial only with its complement, so $\int x^{2}$ is even; at $n=2$ the
  point $|u|=4|c_{2}|$ of $\theta^{2}+\omega$ has $\rho_{2}=23$. In a longer
  run (A) is checked for $n=4$ in every degree, where $HT^{8}$ has
  dimension $12870$, and for $n=5$ up to degree three. This is
  \Cref{prop:p2primeprofile}.

\item \emph{Descent and scalar extension.} The computation models the
  cohomology of an abelian variety of $(F,n)$-Weil type, $F$ a CM field of
  degree $2m$, by alternating $\QQ$-multilinear forms on $H_{1}=F^{N}$, its
  Weil classes by the forms $\Tr_{F/\QQ}(a\det_{F})$ and its
  polarisation by $\Tr_{F/\QQ}(\xi H)$. (A) For $F=\QQ(i)$, $\QQ(\sqrt{-2})$,
  $\QQ(\sqrt{-3})$, $\QQ(\zeta_{5})$ and $\QQ(\zeta_{8})$, and $n=1,2$ (and
  $n=3$ over $\QQ(i)$), the correspondence
  $P(x)=pr_{B*}(x\cdot pr_{Y}^{*}(\eta_{Y}^{2m-2}y'))$ maps the $2m$-dimensional
  space $W_{F}(B\times Y)$ onto $W_{F}(B)$, and it kills $W_{F}(B\times Y)$
  when the Weil class $y'$ of $Y$ is replaced by $\eta_{Y}$. (B) For
  $H_{Y}=\mathrm{diag}(1,-t)$ the form $H_{B}\oplus H_{Y}$ has signature
  $(n+1,n+1)$ and normalised discriminant $t$ times that of $H_{B}$. (C) For
  $\QQ(i)\subset\QQ(\zeta_{8})$, $\QQ(\sqrt{-7})\subset\QQ(\zeta_{7})$ and
  $\QQ(\sqrt{-3})\subset\QQ(\zeta_{9})$, and $n=1,2$, the pull-back
  $j^{*}W_{F}(B_{F})$ is all of $W_{K}(B)$. In a longer run (A) is checked up
  to $n=5$ over $\QQ(i)$, up to $n=4$ over $\QQ(\sqrt{-2})$ and
  $\QQ(\sqrt{-3})$, up to $n=3$ over $\QQ(\sqrt{-5})$, $\QQ(\zeta_{5})$ and
  $\QQ(\zeta_{8})$, and up to $n=2$ over the sextic field $\QQ(\zeta_{7})$, and
  (C) up to $n=3$. This is \Cref{prop:descent}.

\item \emph{Quartic fields, the Mumford fourfold, and the objects of the
  criterion.} Each computation of this item is carried out twice, by
  independent implementations. A fast selection of them performs $58$
  checks; the longer runs are recorded with the programs. Where a space of
  invariants is bounded above by a rank modulo a prime, which can only lower
  a rank, the bound is matched by explicit invariants. (A) Quartic fields:
  for $F=\QQ(\zeta_{5})$ (cyclic, with no imaginary quadratic subfield),
  $\QQ(\zeta_{8})$, $\QQ(\zeta_{12})$ (biquadratic) and
  $\QQ(\sqrt{-(3+\sqrt2)})$ (not Galois), and six hermitian forms in all,
  $\mathfrak{su}(V,H)$ has dimension $30$ and kills $W_{F}$, and its
  invariants have dimension $2$ in degree two and $7$ in degree four, with
  $W_{F}$ meeting the products of divisor classes in $0$; $\phi$ is
  semilinear with $\phi^{2}=(\det H)^{-1}$, it commutes with $\mathfrak{su}(V,H)$
  on $R_{F}$ and the commutant has dimension $8$, and, over $\QQ(\zeta_{5})$
  with a hermitian form that is not diagonal,
  $\phi g=\overline{\det g}\,g\phi$ for an element $g$ of $\mathrm{U}(V,H)$ with $\det g\notin\QQ$;
  \eqref{eq:quarticstar} holds, $\Delta^{*}C_{q}=6c\,\Omega$, and the
  pull-backs of $W_{F}$ along five maps span a space of dimension $20$
  containing $C_{q}$; on the split locus for three fields $\det H$ is a square
  in $F_{0}$ and the polarisation of $X$ gives a Hodge class in $R_{F}$; on
  the scalar-extension locus for three biquadratic fields and four forms the
  pull-backs $g_{f}^{*}W_{K}(C)$ span a space of dimension $10$ containing
  $W_{F}(B_{F})$, and the invariants of $\mathfrak{su}(V_{C},H_{C})$ in degree
  two meet $R_{F}$ in $0$; $\dim HT^{2}=120$, the annihilator of a Weil class
  has dimensions $0$, $16$, $24$, the eight-dimensional tangent space kills
  $\gamma$, and $r(\gamma)=112$ for nine general $p$ and at six elements of
  the Hodge ring; the part of the contraction that no $p$ reaches has rank
  $68$; and on the $25$ monomials and all $300$ pairs of them, with seven ratios, $r$ takes the values $80$,
  $84$, $88$, $104$, $108$, $112$, and $92$, $96$ on pairs, with no value
  below $80$ in about $1800$ further exact trials. This is
  \Cref{prop:quarticquat}, \Cref{prop:quarticnl}(iii),
  \Cref{prop:quarticscalar}, \Cref{prop:quarticcasimir} and
  \Cref{thm:quarticp2prime}. (B) The Mumford fourfold: by peeling highest
  weights and, independently, by Weyl's alternating sum, no
  $\bigwedge^{k}V$ contains a summand $S^{2}V_{i}$, the multiplicity of
  $S^{2}V_{1}$ in $H^{n}(X\times X)$ is $1,6,16,24,16,6,1$ for
  $n=2,4,\dots,14$ and $0$ otherwise, and the copy in $H^{2}$ lies in
  $H^{1}\otimes H^{1}$ with Hodge numbers $(1,7,1)$, as it does in the
  twelve-dimensional unitary model of \Cref{rem:mumfordunitary}; on the split
  model, in the exterior algebras of $V\otimes\QQ^{m}$, the degree-four Hodge
  classes of $X\times X$ have dimensions $1,1,4,1,1$ in the K\"unneth
  components, $(V^{\otimes4})^{G}$ has dimension $8$, the pull-backs along $45$
  integer homomorphisms $X^{4}\to X^{2}$ and the products of pulled-back
  divisor classes span exactly $7$ dimensions and miss $\mathrm{Det}$, six
  times Cayley's hyperdeterminant on the diagonal, and composites of two
  correspondences fill all $8$; $s_{k}=\mathrm{Cas}_{k}+\frac12$, the images of
  $\QQ[S_{2}]$, $\QQ[S_{3}]$, $\QQ[S_{4}]$ have dimensions $2$, $5$, $14$, so
  $\End_{G}(V\otimes V)$ and $\End_{G}(V^{\otimes3})$ have dimensions $8$ and
  $125$; and the irreducible zero-sum multisets of the weights $\{\pm1\}^{3}$
  are the four pairs $\{w,-w\}$ and two sets of four. This is
  \Cref{prop:mumfordwhere}, \Cref{thm:mumfordpowers}, \Cref{rem:mumforddet}
  and \Cref{prop:mumfordcmpowers}; \Cref{cor:mumfordksproducts},
  \Cref{prop:mumfordkslink}, \Cref{prop:mumfordcarrier}, apart from its
  computational part, and \Cref{prop:mumfordsurface} are proved in the text
  and need no computation. (C) Explicit objects on the split member: the
  Hochschild profile of a secant class for $n=2,\dots,5$ and the Euler
  characteristic for $n=2,\dots,6$, at further values of $d$; the
  certificates of \Cref{thm:twistedcert}, with profile $(1,8,18,8,1)$ for seven values of
  $d$ and $(1,12,48,74,\dots)$ for five; the flatness of
  $\ch(E)e^{\ell/2}$, the recursion $\mu_{m+2}=-\mu_{m}/(4d)$ and the Weil
  part over the function field $\QQ(d)$ for $n\le5$, the $B$-field in the
  eight conventions for $\Phi$, and the dimensions $5$ and $9$ of the Hodge
  loci; $r(\ch G)=18$, $48$, $88$ at $n=2,3,4$; the parity numbers of
  \Cref{rem:untwisting}; and the formulas of \Cref{ex:splitsmall} at $n=3$.
  (D) Natural objects at $n=4$: over $\QQ(i)$, exactly, $\dim T=16$,
  $T$ kills $\eta$ and the Weil plane, and $T\lrcorner\ch F$ has rank $6$ for
  the natural objects off $C_{\eta}$ and $0$ on it; modulo a prime,
  $r(\gamma)=56$, $72$, $88$, $104$ for $\rho=0,\dots,3$; the Hodge ring of
  dimension $55$, the fit of $\Phi_{4}$ and the geometry of $C_{\eta}$,
  $l_{W}$ and $L_{\pm}$ in \Cref{setup:lagquadric}; the relation of
  \Cref{ex:fourteen} in both sign conventions, its minimality among the
  $12870$ eight-sets, the configurations of \Cref{thm:minimalsupport} (eight
  points on one conic give a pure relation, and the splittings $7$, $4+4$,
  $5+3$, $6+2$, $6+1+1$, $7+1$, $3+3+2$ give none), the lattice spanned by
  the classes of the subvarieties, of index $2612736000$, and the intersection
  numbers of the sum of \Cref{ex:fourteen}; and the agreement with the
  exhaustive search of item (XIII) at $n=2$, where the Weil tuples all use
  four points off $C_{\eta}$ and the pure ones are those on a conic through
  $l_{W}$. \Cref{lem:p2primesummands}, \Cref{cor:naturalsums},
  \Cref{prop:secantprofile} and \Cref{prop:orlovfour} are proved in the text
  and checked here in examples. (E) Secant classes over quartic and sextic
  fields, and four further statements, in their smallest cases: for a
  quartic CM field, the secant spaces $S(t,q)$ are rational of dimension $4$, with Hochschild profile $(1,8,2\rho+4N_{w})$ and
  the Euler characteristic of \Cref{prop:quarticsecant}(iii), in $49$ cases;
  the expansion in the proof of \Cref{lem:quarticlattice}, with $\delta$ and
  the entries of $V$ as symbols, and a direct computation of the integral
  classes of $\mathcal{R}$, which form the congruence lattice $c\equiv g$
  modulo the different for the discriminants $5$, $8$, $13$, $17$ and $40$ in
  the product model $S\otimes M$ and for $12$, $24$, $28$, $40$ and $56$ on
  the lattice $\cO^{2}\oplus(\mathfrak d^{-1})^{2}$; the arithmetic of
  \Cref{prop:quarticother}, namely $\Nm(x+4y)\equiv\Nm(x)$ modulo $4$ with
  $D$ a symbol, the parity at the ramified prime above $2$, and the three
  conditions on $m$ for every squarefree $D<200000$, which leave only $D=2$
  and $D=5$; an exact scan of $1841$ complex-secant planes for
  $\QQ(\sqrt{10})$, which finds no admissible character; the invariance of
  the twisted character of an Orlov product of such objects under the Lie
  algebra of $G_{F}$, and its nonzero Weil part, for the seven real
  quadratic fields $\QQ(\sqrt D)$, $D=2,3,5,7,10,13,17$, with at least two
  values of $q$ each; the formula $r(\kappa)=r^{2}(v_{1})+64+r^{2}(v_{2})$
  of \Cref{thm:quarticobstruction}(i) on eight pairs, from an exact annihilator
  over $\QQ(\sqrt{-1})$, with the values $100$, $102$, $104$ and, for
  $F=\QQ(\sqrt5,\sqrt{-1})$ and a class with $N_{w}=2$, $88$ and $94$; for
  \Cref{prop:sexticcount}, an independent computation of the full profiles
  $r^{0},\dots,r^{6}$ of the secant classes of a sextic field, for every
  support closed under $\varepsilon\mapsto-\varepsilon$ and three tensors of
  smaller flattening rank, exactly over $\ZZ$ at two random points, which are
  palindromic and agree with item (LVI), with the value of $\chi$ over
  $\QQ(\sqrt{-1})$ and the quartic profile as controls; and the identities
  used in the proof of \Cref{prop:quarticother}, on random integral classes.
  Further, for \Cref{prop:flatall}, the four identities at $n=1$ and $n=2$
  with $d$ a symbol; for \Cref{prop:orlovequality}, the Euler characteristic
  of a secant class and the numbers $R_{n}$ for $n\le12$; for
  \Cref{prop:divisortemplate}, the profiles of part (i) for $5\le n\le24$ on
  $4320$ secant characters, the series identities of parts (ii) and (iv) with
  $b$, $d$ and $r$ as symbols, the counts of part (iii) for $n\le39$, and
  the inequality closing part (iv); for
  \Cref{prop:orthpowers}, the dimensions of the $\SO(t)$- and $\mathrm{O}(t)$-invariants of $V^{\otimes j}$,
  equal for $j<t$ and differing by one at $j=t$, for $t\le8$, with the
  determinant invariant and odd under reflections; and for
  \Cref{prop:f3primestrength}, the closure graph with the implication
  $\textup{(F3$'$)}\Rightarrow\textup{(M)}$ added, which has the same five
  minimal sufficient sets as the rule set of item (XXXIII).

\item \emph{The count in degree six.} The computation builds $H^{*}(X,\CC)$
  for a principally polarised abelian sixfold with real multiplication by a
  cubic field as the exterior algebra on twelve generators in three blocks of
  four, with $\theta_{j}=x_{j1}y_{j1}+x_{j2}y_{j2}$, and $HT^{k}(X)$ as
  $\bigoplus_{a+b=k}\bigwedge^{a}H^{0,1}\otimes\bigwedge^{b}T$ acting by
  $\alpha\wedge(t\lrcorner\,\cdot\,)$. (A) $\theta_{j}^{3}=0$,
  $\int\theta^{6}=720$, $\int\theta_{1}^{2}\theta_{2}^{2}\theta_{3}^{2}=8$ and
  $\int\exp(\sum c_{j}\theta_{j})=(c_{1}c_{2}c_{3})^{2}$. (B) On one factor,
  $HT^{k}_{j}$ kills $e^{\lambda\theta_{j}}$ for $k\ge3$, $HT^{2}_{j}$ maps
  into the line of $\omega_{j}$, and the coefficient pairs at
  $\lambda_{j}(\pm1)$ span $\CC^{2}$. (C) Exactly over $\QQ(\sqrt{-1})$, for
  $q=1$, $t\in\{0,2\}$ and four coefficient patterns (generic, supported on
  the two CM types induced from the imaginary quadratic subfield, a rank one
  tensor, and a Galois orbit of six types), the rank of contraction from
  $HT^{1}$ is $12$ and the ranks from $HT^{2}$ and $HT^{3}$ are $(54,112)$,
  $(30,40)$, $(51,88)$, $(54,96)$, equal to the formulas of \Cref{prop:sexticcount}(i) evaluated on
  $N_{w}$, the $A_{jj'}$, the $\rho_{j}$ and the slices $M^{(j;j')}_{\sigma'}$.
  (D) Exactly, $\chi(v,v)=-64\Nm(q)\sum|w_{\varepsilon}|^{2}$ for the four
  patterns, and $+16\Nm(q)\sum|w_{\varepsilon}|^{2}$ for two places. (E) Modulo
  a prime $p\equiv1$ modulo $4$ near $10^{6}$ at which the cubic field splits
  and the $s_{j}$ exist, for $\QQ(\zeta_{7})^{+}$ with $q=2+\alpha$,
  $\QQ(\zeta_{9})^{+}$ with $q=3-\alpha$ and the non-Galois cubic field of
  discriminant $229$ with $q=3+\alpha$, $\alpha$ a generator and $t=\alpha$,
  the ranks modulo $p$, which bound the ranks below, equal the values of the
  formulas, which bound them above, and the identity of (D) holds modulo
  $p$. (F) The bookkeeping of \Cref{prop:sexticcount}(iii): the bound
  $\chi\le-(r^{3}-2r^{2}+22)$ is $-26$, $-2$, $-8$, $-10$ for the four
  patterns, and $r(\kappa)\le252<264$. This is \Cref{prop:sexticcount}.

\item \emph{Integral flat characters in degree six.} (A) On a complex torus
  of dimension six, $\int v^{\vee}v$ is even for forty random integral
  classes of even degree. (B) Over the $255$ nonempty supports of a
  $2\times2\times2$ coefficient tensor and all ranks of its six slices and
  of its three flattenings allowed by the constraints in the proof of
  \Cref{prop:sexticintegral}(ii), $r^{3}-2r^{2}$ is at most $4$, and equal to
  $4$ only at $(r^{2},r^{3})=(54,112)$, which repeats the case analysis of
  that proof. (C, D) For the sixteen totally real
  cubic fields $\QQ(\zeta_{7})^{+}$, $\QQ(\zeta_{9})^{+}$ and the fields
  generated by a root $\alpha$ of $x^{3}+ax^{2}+bx+c$ with $|a|\le1$,
  $|b|,|c|\le7$, irreducible with squarefree discriminant (the discriminants
  $229$, $257$, $321$, $469$, $473$, $697$, $761$, $785$, $985$, $993$,
  $1129$, $1229$, $1345$ and $1509$, for which $\cO=\ZZ[\alpha]$), and for
  $q=1$ and $q=k+\alpha$, the lattice of integral points of the flat secant
  space $S(0,q)$ in the model
  $\cO e_{1}\oplus\cO e_{2}\oplus\mathfrak d^{-1}f_{1}\oplus\mathfrak d^{-1}f_{2}$:
  the coordinates of the eight parameter classes in the monomial basis of
  $\bigwedge^{\mathrm{ev}}H^{1}(X,\ZZ)$, of dimension $2048$, are computed
  modulo three primes near $2^{40}$ that split completely, reconstructed as
  rationals and confirmed modulo a fourth; the lattice is the dual of the
  span of the rows; the Gram matrix of $-\chi/8$ from
  \Cref{prop:sexticintegral}(iii) is integral and agrees with the direct
  integral of $v^{\vee}v$ modulo a prime near $2^{45}$; a Fincke--Pohst
  enumeration finds no vector with $-\chi\le24$ in any of the thirty-two
  cases, and for $q=1$ exactly the four vectors $\pm\operatorname{Re}
  e^{\sqrt{-1}\,\theta}$, $\pm\operatorname{Im}e^{\sqrt{-1}\,\theta}$ with
  $-\chi=32$, each of shape $N_{w}=2$, the shape being computed exactly
  modulo a prime near $10^{6}$ at which the $s_{j}$ exist. This is
  \Cref{prop:sexticintegral}(i), (ii) and (iv); the lattice minima of (C, D)
  are not used in the paper.

\item \emph{The $F$-Weil part of Orlov products in degree six.} The
  computation works in the exterior algebra on the twenty-four generators of
  $H^{1}(X\times\widehat X)$, three factors of eight, with $\sqrt{-q}$ acting
  on the $j$-th factor as in \eqref{eq:Mdef} with $\beta=\theta_{j}$ and
  $d=d_{j}$, and uses the exact Fourier--Mukai transform of item (LV) and a
  version pruned by output degree, which agrees with it in degrees at most
  four. (A) The vectors $x_{i}+\sigma(s_{j}/d_{j})Bx_{i}$ are eigenvectors of
  $\sqrt{-q}$ with eigenvalue $\sigma s_{j}$, and their wedge product is the
  class $\omega_{j,\sigma}$ of \eqref{eq:sexticomega}; on one factor,
  $\Psi(e^{\sigma s\theta}\boxtimes e^{\sigma s\theta})=-4d\,
  e^{\sigma(s/2d)\eta}$ with $\eta=d\theta+\widehat\theta$ and
  $\Psi(e^{\sigma s\theta}\boxtimes e^{-\sigma s\theta})
  =\frac12(\gamma-\sigma s\ell)^{2}$ for $d=1,4,\frac14$ and $\sigma=\pm1$;
  and on the twelvefold, exactly over $\QQ(\sqrt{-1})$ for
  $(d_{1},d_{2},d_{3})=(1,1,1)$, $(1,4,9)$ and $(4,\frac14,\frac94)$, the
  identities being formal in the $d_{j}$, and three pairs each, one of them
  diagonal, the degree-four part of $\kappa$ lies in the span of the
  $\eta_{i}\eta_{k}$ and the $\omega_{j,\sigma}$, with the coefficients of
  \Cref{lem:sexticweil}(i). (B) For $q=1$, the part of degree at most four of
  \Cref{lem:sexticweil}(iii) on six pairs in the plane of
  $\operatorname{Re}e^{\sqrt{-1}\,\theta}$ and
  $\operatorname{Im}e^{\sqrt{-1}\,\theta}$; no $F$-Weil part for two classes
  supported on one pair $\{\varepsilon_{0},-\varepsilon_{0}\}$; and
  $\chi(v(-5,0,-2,0))=-296$, the pair $(\operatorname{Re}
  e^{\sqrt{-1}\,\theta},v(-5,0,-2,0))$ giving rank $-8$ and all six
  $\kappa_{j,\sigma}$ equal to $7$, the diagonal pair rank $-296$ and all
  six equal to $-21$. (C) The identity $\sum_{\varepsilon_{j}=\sigma}
  w_{\varepsilon}w_{\varepsilon^{(j)}}=\frac1{16}\sum_{b}D_{b}^{2}/
  \prod_{i\in b}d_{i}$ of the proof of \Cref{lem:sexticweil}(ii) on random
  classes; the expansion with $u^{2}=-1/q_{j}$, symbolically at each place;
  the formulas for $R$ and $I$ in $F_{0}$ against the conjugate formulas
  modulo a prime that splits $F_{0}$, in the forty-eight cases below; exactly
  over $\QQ(\alpha)$ on the twelvefold, for $\QQ(\zeta_{7})^{+}$ and
  $\QQ(\zeta_{9})^{+}$, $q=2+\alpha$ and $3+\alpha$ and five classes, two of
  them with $I\ne0$ and one with $\Omega=0$, the formula
  $\kappa_{j,\sigma}(v,v)=-\Nm(q)\tau_{j}(q)\tau_{j,\sigma}(\Omega(v))$; and
  the two values of $\Omega$ at the minima of (E). (D)
  $\beta^{2}=2+\alpha$ with $\Nm(\beta)=-1$ and $1$ for the two cyclic
  fields, and the integral classes of $S(0,2+\alpha)$ with $-\chi\le32$,
  which are $\pm\operatorname{Re}e^{\sqrt{-1}\,\theta_{\beta}}$ and
  $\pm\operatorname{Im}e^{\sqrt{-1}\,\theta_{\beta}}$, with $\Omega=0$. (E)
  For the sixteen cubic fields of item (LVII) and $q\in\{1,k+\alpha,
  k+1+\alpha\}$, forty-eight cases, the lattice of integral points of
  $S(0,q)$ as in item (LVII), reduced by LLL and enumerated exactly by
  Fincke--Pohst with a bound doubled until a class with $\Omega\ne0$
  appears: the least $-\chi$ with $\Omega\ne0$ is $192$, attained only at
  $\pm v(-1,0,2-\alpha^{2},0)$ over $\QQ(\zeta_{7})^{+}$ with $q=3+\alpha$;
  the lattice minima $192$ and $256$ for $q=3+\alpha$ over the two cyclic
  fields, each attained by one pair $\pm v$; for $q=1$ over the two cyclic
  fields, all classes with $-\chi$ below $296$, respectively $488$, in
  $\ZZ\operatorname{Re}e^{\sqrt{-1}\,\theta}\oplus\ZZ\operatorname{Im}
  e^{\sqrt{-1}\,\theta}$, and sixteen classes at those values, among them
  $v(-5,0,-2,0)$ and $v(-7,0,-2,0)$, with the same least values for
  $q=2+\alpha$; the least values between $192$ and $7080$ in each case, and
  between $296$ and $3872$ for $q=1$. The shapes $N_{w}=8$,
  $(r^{2},r^{3})=(54,112)$ are computed modulo a prime near $10^{6}$ at which
  the $s_{j}$ exist; ranks modulo a prime bound the ranks below, and these
  are the largest values. (F) The imaginary quadratic subfields of
  $F_{0}(\sqrt{-q})$ in the forty-eight cases, from the factorisation over
  $\QQ$ of $\prod_{j}(x^{2}-D\tau_{j}(q))$ with $D$ the squarefree part of
  $\Nm(q)$: exactly twenty cases, with
  the lattice minima $288$ and $4608$, attained with $\Omega=0$, and the
  least values $728$ and $6504$ with $\Omega\ne0$ in the two cases with
  $\QQ(\sqrt{-3})$. Parts (A) to (C) are \Cref{lem:sexticweil}; the lattice
  minima of (D) to (F) are not used in the paper.

\item \emph{Kernels on $X\times X$ for a quartic field.} The computation
  works over $\RR$ place by place, as in the proof of
  \Cref{thm:quarticobstruction}(iii). At a real place $\tau_{j}$, with
  $d=\tau_{j}(q)$, it works in the exterior algebra on the eight generators
  of the part of $H^{1}(X\times\widehat X)$ at $\tau_{j}$, with the classes
  $\theta_{j}$, $\widehat\theta_{j}$, $\ell_{j}$, $\eta_{j}$, $\gamma_{j}$,
  the Lie algebra of $SU_{j}$ and an exact Fourier--Mukai transform; a class
  on $A$ is a sum of tensor products over the two places, split exactly along
  the seven invariant classes at each place. (A) At one place: the transform
  of the classes of $X\times\{0\}$, of the diagonal, of the antidiagonal and
  of six graphs; $\Phi_{j}(P_{j})$ is minus the factor of the class of
  $X\times\{0\}$ in $A$, and is not flat; the formula
  $\ch\Phi_{j}(T_{j}-P_{j})=-\frac14(\eta_{j}^{2}+\gamma_{j}^{2}-d\ell_{j}^{2})$
  for the mixed kernel $T_{j}-P_{j}$ at $d=1,2,3,\frac52$, both sides being
  of degree at most two in $d$, and over $\QQ(\sqrt5,\sqrt{-1})$ at
  $d=\frac12(5+\sqrt5)$; the flatness of this class, its pure degree four,
  its nonzero Weil part, and its Hodge type at $X\times\widehat X$; the
  eigenvalues $\pm4s_{j}$ of $\sqrt{-q}$ on $(\gamma_{j}\mp s_{j}\ell_{j})^{2}$;
  the non-flatness of its twists by $e^{\pm\ell/2}$; the ratio
  $\int\eta_{j}^{4}:\int(\gamma_{j}^{2}-d\ell_{j}^{2})^{2}=3:4$; and its rank
  of contraction $23$, with annihilator of dimension $5$ in $H^{1}(T)$,
  against $24$ and $4$ for a general invariant class of degree four. (B) For
  $\QQ(\zeta_{5})$ with $q=2+\frac12(1+\sqrt5)$, $\QQ(\sqrt5,\sqrt{-1})$ with
  $q=1$, $\QQ(\sqrt2)(\sqrt{-q})$ with $q=2+\sqrt2$, and
  $\QQ(\sqrt{13})(\sqrt{-q})$ with $q=2+\frac12(1+\sqrt{13})$, the class
  $\ch\Phi(Z'')$ has pure degree four, $r^{1}=16$ and $r=94$, and its
  annihilator of dimension $26$ consists of the sixteen bivectors across the
  two places and a $5$-dimensional subspace of $H^{1}(T)$ at each place. (C)
  For five pairs $(F_{0},q)$, $G''$ is invariant under
  $\mathrm{Gal}(F_{0}/\QQ)$, and the classes $(\beta,\alpha)_{*}(c)$ with
  $\beta,\alpha\in\cO$ span all $81$ rational classes of graph type at both
  places; for $\QQ(\zeta_{5})$, $G''$ is an explicit rational combination of
  $36$ classes $(\beta,1)_{*}(c)$ on $21$ graphs; and
  $\frac12(\theta_{1}^{2}+\theta_{2}^{2})
  =\frac1{10}(3\theta^{2}-2\theta\theta_{R}+2\theta_{R}^{2})$ for
  $F_{0}=\QQ(\sqrt5)$, with $\theta_{R}$ the class of the real multiplication
  by $\frac12(1+\sqrt5)$. (D) For $F_{0}=\QQ(\sqrt5)$, $q=\frac12(5+\sqrt5)$,
  and $F_{0}=\QQ(\sqrt2)$, $q=2+\sqrt2$, and six pairs of characters in
  $S(0,q)$ each, five of the twelve of rank zero: $\kappa=\ch(E)e^{\ell/2}$ is
  flat and $\ch(E)$ is not; the relation between $\kappa_{(4,8)}$ and
  $\kappa_{(4,0)}$ of \Cref{rem:quarticsecantshape} and its symmetric form;
  and, for the five pairs in each field where $\kappa_{(4,0)}$ has a Weil
  part, components of $\kappa$ outside $\CC[\eta_{1},\eta_{2}]+W_{F}\otimes\CC$.
  (E) At one place: the classes $(b,a)_{*}(c)$ with $c$ of even degree span a
  space $L_{j}$ of dimension $24$, and $\Phi_{j}(L_{j})$ meets the invariants
  exactly in $\mathrm{span}(1,\mathrm{pt}_{j})$ for $d=1,3,\frac25$; the
  $2$-forms $B$ with $B\wedge\Phi_{j}(L_{j})\subseteq\Phi_{j}(L_{j})$ form a
  space of dimension $7$, spanned by the $2$-forms pulled back from $X$ and
  by $\ell_{j}$, not containing $\widehat\theta_{j}$; $e^{B}\Phi_{j}(L_{j})$
  meets the invariants exactly in $e^{g\eta_{j}}\mathrm{span}(1,\mathrm{pt}_{j})$
  for four $B=f\theta_{j}+g\widehat\theta_{j}+h\ell_{j}$, one with
  $h=-\frac12$; and for the graph of $\frac12(1+\sqrt5)$ in $\QQ(\zeta_{5})$,
  $L=\cO$ and $L=p_{1}^{*}\cO(\theta)$, and $t\in\{0,\frac12,1\}$, the class
  $\ch\Phi(I_{\Gamma}\otimes L)e^{t\ell}$ is not flat. (F) At one place: for
  $d=1,3,\frac25$ the commutant of the Lie algebra of $SU_{j}$ on the part
  of $H^{1}(A)$ at $\tau_{j}$ is spanned by $1$ and $\sqrt{-q}$, so that part
  is irreducible, and the invariants have dimensions $1,1,3,1,1$ in degrees
  $0,2,4,6,8$; the transforms of line bundles on seven graphs are pure
  spinors, with annihilator of dimension $8$, also after the twist by
  $e^{\ell/2}$; and for $d=1$ the flat classes $\ch\Phi_{j}(T_{j}-P_{j})$ and
  $\gamma_{j}^{2}-d\ell_{j}^{2}$ are not pure, while
  $(\gamma_{j}-s_{j}\ell_{j})^{2}$ is pure over $\CC$. (G) $\int_{X}x^{2}$ is
  even for forty random integral classes $x$ of degree four on a torus of
  dimension four; for $\chi=4$ and $(\rho,N_{w})=(2,4)$ the conditions (a) to
  (c) of \Cref{rem:quarticnegext} leave only $D=3$, $m=1$ among the squarefree
  $D<200000$; and the resulting bounds on $\nu_{i}$. (H) $\varphi(N)=4$
  exactly for $N=5,8,10,12$ among $N<200$; $3g-3<\frac12\varphi(N)(g-1)^{2}$
  whenever $\varphi(N)\ge4$ and $g\ge3$, for $N<200$ and $g<60$; and a
  hyperbolic hermitian form of rank $2n$ has determinant $(-1)^{n}$. This is
  \Cref{rem:quarticsecantshape}, \Cref{rem:quarticnegext},
  \Cref{prop:quarticexclusions} and the numbers in \Cref{rem:quarticprym};
  the kernels of (A) to (E) are not used in the paper, and the algebraicity
  of the Weil classes of the Prym varieties is Schoen's theorem and is not
  recomputed.

\item \emph{The motivic group of a Mumford fourfold, the known classes, and
  the Kuga--Satake base points.} The computation works in the split model
  $V=\QQ^{2}\otimes\QQ^{2}\otimes\QQ^{2}$ with
  $\psi=\epsilon\otimes\epsilon\otimes\epsilon$ and $G=\SL_{2}^{3}$ acting
  factorwise. (A) $\mathfrak{sp}(V,\psi)$ is the sum of
  $\operatorname{Lie}G$ and the span $P$ of the $27$ tensors
  $x_{1}\otimes x_{2}\otimes x_{3}$ with $x_{k}\in\mathfrak{sl}_{2}$; $P$ is
  stable under $\operatorname{Lie}G$, the commutant of its action has
  dimension one, and by weights $S^{2}V$ is multiplicity free, with summands of
  highest weights $(2,2,2)$, $(2,0,0)$, $(0,2,0)$, $(0,0,2)$. (B) The six
  permutations of the factors lie in $\Sp(V,\psi)$, normalise
  $\operatorname{Lie}G$ and permute its factors accordingly; the commutant of
  $\operatorname{Lie}G$ in $\End V$ is $\QQ$; $-1$ is the image of
  $(-1,1,1)$; and the subgroups of $S_{3}$ stable under conjugation by $A_{3}$
  are $1$, $A_{3}$ and $S_{3}$. (C) The projectors $\pi_{0}$, $\pi_{12}$,
  $\pi_{13}$, $\pi_{23}$ of \Cref{setup:omegaa} have ranks $1$, $9$, $9$,
  $9$, and the cyclic permutation $\zeta$ carries $\pi_{12}$ to $\pi_{23}$,
  $\pi_{23}$ to $\pi_{13}$ and $\pi_{13}$ to $\pi_{12}$ and fixes $\pi_{0}$.
  (D) $r_{1}=s_{1}+s_{2}+s_{3}$ commutes with $\operatorname{Lie}G$ and with
  the six permutations, acts on $S^{2}V$ by $3$ on the summand of dimension
  $27$ and by $-1$ on the other three, and does not commute with
  $\mathfrak{sp}(V,\psi)$; Cayley's hyperdeterminant, a sum of $12$
  monomials, is killed by $\mathfrak{sl}_{2}^{3}$, fixed by the six
  permutations and not killed by $\mathfrak{sp}(V,\psi)$; and by Weyl's
  alternating sum $\dim(S^{4}V)^{G}=1$ and $\dim(S^{4}V)^{\Sp(V,\psi)}=0$.
  (E) For $m=1,2,3$ the products over the three factors of the non-crossing
  pairings form a basis of $(V^{\otimes2m})^{G}$, of dimension $1$, $8$,
  $125$ as Weyl's sum gives, which $\zeta$ and a transposition of the factors
  permute; counting orbits, $\dim(V^{\otimes2m})^{H}$ for $H=G$, $G\cdot
  A_{3}$, $N$, $\Sp(V,\psi)$ is $8$, $4$, $4$, $3$ for $m=2$ and $125$, $45$,
  $35$, $15$ for $m=3$, the last column being spanned by the $(2m-1)!!$
  pairings of $\psi$. (F) $(V\otimes V)^{\operatorname{Lie}G}=\QQ\psi$; the
  divisor classes of $X\times X$ are $\mathfrak{sp}(V,\psi)$-invariant; the
  Weil classes of $X\times X$ for $\QQ(\sqrt{-1})$ acting through
  $\mathrm{M}_{2}(\QQ)$, the Weil classes of $X\times X\times B$, $B$ an
  abelian surface, for $\QQ(\sqrt{-1})$ acting diagonally, and the orientation
  class of $X\times X$ are $\mathfrak{sl}(V)$-invariant. (G) For
  $F=\QQ(c)$, $c=2\cos(2\pi/7)$, and $D=(-1,c)_{F}$, which is definite at
  exactly the two real places where $c<0$ and has split corestriction since
  $\Nm(c)=1$, the space $(T,q_{1})$ has signature $(2,7)$ and contains two
  explicit hyperbolic planes, with negative definite complement of rank $5$,
  and a subspace $W$ of signature $(2,3)$ and Witt index two with negative
  definite complement of rank $4$; the lattices $U\oplus\langle2a\rangle
  \oplus\langle-2b\rangle$ and $U^{2}\oplus\langle-2m\rangle$ embed
  primitively in $U^{3}$ for three values each; for an abelian surface with
  $\NS=\langle h\rangle$, $h^{2}=2d$, the transcendental lattice contains
  $U(-1)\oplus U\oplus\langle-2d\rangle$ for $d=1,2,3,5$;
  $\langle6,-2,-2\rangle$ is anisotropic over $\QQ_{3}$, by the Hilbert
  symbol and by descent, so $U\oplus\langle6\rangle\oplus\langle-2\rangle
  \oplus\langle-2\rangle$ has Witt index one; and this lattice, of signature
  $(2,3)$, embeds primitively in $U^{3}\oplus E_{8}(-1)^{2}$. This is the
  computational part of \Cref{prop:mumfordmotivic}, \Cref{prop:mumfordformal},
  \Cref{rem:mumfordks} and \Cref{rem:mumfordpowersopen}.

\item \emph{Hodge classes on K3 surfaces and on varieties of K3$^{[n]}$
  type.} (A) The Betti numbers of $S^{[2]}$ for a K3 surface, from the
  decomposition of de Cataldo and Migliorini and from G\"ottsche's product
  formula, are $(1,0,23,0,276,0,23,0,1)$, so $b_{4}=276=\dim
  \mathrm{Sym}^{2}H^{2}$ and the Euler number is $324$. (B) The form on
  $\mathrm{Sym}^{2}V$ given by the polarised Fujiki relation has
  determinant $\det(G)^{m+1}2^{m-1}(m+2)$ in the monomial basis, on fifteen
  random forms of ranks $2$ to $6$ and, for the Beauville--Bogomolov lattice
  $U^{3}\oplus E_{8}(-1)^{2}\oplus\langle-2\rangle$, on the $276\times276$
  matrix, where it is $2^{46}\cdot25$. (C) $\int q^{\vee}xy=25\,q(x,y)$ and
  $\int c_{2}xy=30\,q(x,y)$ for $c_{2}=\frac65q^{\vee}$ on all pairs of
  basis vectors, $\int c_{2}^{2}=828$, and with $c_{4}=324$,
  $\operatorname{Todd}_{4}=(3\cdot828-324)/720=3=\chi(\cO_{X})$, so that
  Riemann--Roch returns $\chi(L)=\binom{q(L)/2+3}{2}$. (D) With
  $\NS=\langle h,e\rangle$, $q(h)=2$, $e$ a root, and $T=\NS^{\perp}$ of
  dimension $21$: $\int c_{\varphi}xy=\Tr(\varphi)q(x,y)+2q(\varphi x,y)$
  for a random self-adjoint $\varphi$, $x\in T$ and $y\in H^{2}$; the
  polarised Fujiki formula for $\int xyh^{2n-2}$ at $n=2,3,4$; for a random
  rational isometry $u$ of $T$, a Cayley transform, and $n=2,3$, the
  restriction to the diagonal of $\sum_{i}b_{i}\otimes(\id_{\NS}\oplus
  u)(a_{i})$ is $\kappa^{-1}c_{(u+u^{*})/2}$ plus a class of
  $\mathrm{Sym}^{2}\NS$; and at $n=2$ the correspondence
  $x\mapsto(L^{2})^{-1}(c_{\varphi}\cup x)$ acts on $T$ as
  $(\Tr(\varphi)+2\varphi)/q(h)$. (E) For every $t\le21$ the least $n$ such
  that a partition of $n$ has $t$ distinct part sizes is $t(t+1)/2$, over all
  partitions of $n\le240$ by a dynamic programme and by listing those of
  $n\le30$; for $t=21$ it is $231$. (F) For $W=\QQ\oplus T$, $T=\QQ^{t}$, and
  a group permuting the factors of $W^{\otimes l}$ within blocks, the
  $\SO(t)$-invariants fixed by the group exceed the $\mathrm{O}(t)$-invariants
  fixed by it exactly when there are at least $t$ blocks, by exact integer
  ranks in twelve cases with $t=3,4$ and $l\le4$; and the average of
  $\det(T)$ on three of four
  factors vanishes exactly when two of them lie in one block. (G) For ten CM
  fields of degree $2$ to $10$, among them $\QQ(\zeta_{m})$ for
  $m=5,7,8,9,11,12,15$, the elements $(1+sy)(1-sy)^{-1}$ with $\bar y=-y$ have
  norm one and span $E$, and their symmetric parts span $E_{0}$. (H) For the
  generalised Kummer varieties $K_{m-1}(A)$, $m=2,\dots,6$, sector by sector
  with the action of $A[m]$: the Euler numbers $m^{3}\sigma(m)=24$, $108$,
  $448$, $750$, $2592$, the known Betti numbers for $m=3,4$, the invariant
  part equal to the Poincar\'e polynomial of $A^{[m]}$ divided by
  $(1+z)^{4}$, the rest equal to $15$, $80$, $(15,345,15)$ in degrees
  $(4,6,8)$ and $624$ for $m=2,\dots,5$, and carried only by sectors that are
  finite sets of points exactly when $m$ is prime. (I) The span of the
  elements with rational square, the compositum of the quadratic subfields,
  is $\QQ$ for $\QQ(\zeta_{7})^{+}$, $\QQ(\zeta_{9})^{+}$ and
  $\QQ(\zeta_{11})^{+}$, $\QQ(\sqrt2)$ for $\QQ(\sqrt{2+\sqrt2})$,
  $\QQ(\sqrt5)$ for $\QQ(\zeta_{15})^{+}$, and all of $\QQ(\sqrt2,\sqrt3)$.
  This is the computational part of \Cref{prop:k3powers},
  \Cref{prop:k3ntype}, the remark after \Cref{prop:aclosure} and
  \Cref{rem:f3primefrontier}.

\item \emph{The reach of the zero-cycle arguments.} (A) The graded pieces
  of the Jacobian ring of a smooth hypersurface of degree $d$ in
  $\PP^{n+1}$, computed as the coefficients of
  $(1+s+\dots+s^{d-2})^{n+2}$ and by inclusion and exclusion over the $n+2$
  generators of degree $d-1$, agree for $n\le8$, $2\le d\le8$ and every
  degree, are symmetric about $(n+2)(d-2)$ and have total dimension
  $(d-1)^{n+2}$. (B) By Griffiths' description, the primitive Hodge numbers
  $h^{4,0},h^{3,1},h^{2,2},h^{1,3},h^{0,4}$ of a smooth sextic fourfold are
  $1$, $426$, $1751$, $426$, $1$. (C) The Euler number obtained from these
  Hodge numbers equals $((1-d)^{n+2}-1)/d+n+2$ and the degree of the top
  Chern class, $d$ times the coefficient of $h^{n}$ in
  $(1+h)^{n+2}/(1+dh)$, for $n\le8$ and $2\le d\le8$; for the sextic
  fourfold it is $2610$, with $b_{4}=2606$. (D) $h^{4,0}$ of a
  hypersurface fourfold of degree $3,\dots,7$ is $0$, $0$, $0$, $1$, $6$,
  equal to $h^{0}(\PP^{5},\cO(d-6))$, and vanishes exactly for the Fano
  degrees $d\le5$. (E) For a smooth sextic fourfold $Y$ in a hyperplane of
  $\PP^{6}$, the blow-up formula gives $\mathrm{Bl}_{Y}\PP^{6}$ the middle
  Hodge numbers $h^{6,0}=0$, $h^{5,1}=1$, $h^{4,2}=426$, $h^{3,3}=1753$,
  $h^{p,0}=0$ for $p>0$, only even cohomology, and total Betti number
  $2617=e(\PP^{6})+e(Y)=e(\PP^{6})+e(E)-e(Y)$. (F) In the model
  $H^{*}(Y)[\xi]/(\xi^{2}+c_{1}\xi+c_{2})$ of a $\PP^{1}$-bundle over a base
  with $H^{*}(Y)=\QQ[y]/(y^{5})$, for twenty random rational $c_{1}$,
  $c_{2}$, with $\rho_{*}$ the coefficient of $\xi$ and $j^{*}j_{*}$
  multiplication by $-\xi$, $-\rho_{*}j^{*}j_{*}\rho^{*}a=a$ and
  $\rho_{*}\rho^{*}a=0$ in every degree. (G) A class of degree $2p$ on an
  $n$-fold needs an input beyond \Cref{prop:f3prime}(iv) and hard Lefschetz
  only for $2\le p\le n/2$, so only $p=2$ for $n\le5$; on a blow-up of a
  sixfold along a smooth centre $C$ the summands of $H^{2p}$ with $p\le3$
  need one only in degree four on a fourfold $C$; and the decomposition of
  the diagonal of a uniruled fivefold needs one only in degree four on a
  fourfold. This is the computational part of \Cref{prop:f3primefibration}
  and \Cref{rem:f3primesharp}.

\item \emph{The Wirtinger bound on a Mumford square.} The computation works
  in the split model of item (XLIV), with $\theta=\psi$ the class of the
  principal polarisation of $X$, $\theta_{Y}=\theta\otimes1+1\otimes\theta$
  on $Y=X\times X$, and the volume forms normalised by $\int_{X}\theta^{4}=4!$
  and $\int_{Y}\theta_{Y}^{8}=8!$. (A) $\psi$ is nondegenerate,
  $\theta^{4}=24\,\mathrm{vol}$ and $\theta_{Y}^{8}=8!\,\mathrm{vol}$.
  (B) Every element of $U_{12}$, $U_{13}$, $U_{23}$ is primitive,
  $x\theta^{3}=0$, while $\theta^{4}=24\,\mathrm{vol}$. (C) The pairings
  $\int_{Y}\pi_{ij}\theta_{Y}^{6}$ vanish for the three nontrivial summands and
  $\int_{Y}\pi_{0}\theta_{Y}^{6}=2880$, where
  $\pi_{0}=\frac14\,\theta\otimes\theta$ in this model; so
  $\int_{Y}\omega_{a}\theta_{Y}^{6}=2880\,a_{0}$, checked on five random
  integer vectors $a$, and $L_{\theta_{Y}}(\omega_{a})=4a_{0}$. (D)
  $L_{\theta_{Y}}(\theta_{Y}^{2})=8!/6!=56$ and
  $L_{\theta_{Y}}(4\theta_{Y}^{2})=224$. That is \Cref{ex:mumfordmass}.

\item \emph{The Weil structure at a CM point and the cycles on the square.}
  In the model of item (XLIV) the Hodge group at a CM point is the diagonal
  torus of the weight basis, and the checks are in exact rational arithmetic
  in the exterior algebra of $H^{1}(X\times X)$. (A) The monomials of weight
  zero number $16$ in degree two and $132$ in degree four; the products of
  the $16$ divisor classes span $100$ dimensions and those of the three
  $\Sp$-invariant divisor classes $6$; the zero-sum four-sets of weights are
  the six unions of two opposite pairs and the two tetrahedra $T_{+}$,
  $T_{-}=-T_{+}$, each meeting every opposite pair once; the $32$ monomials
  of weight zero that are not products of two divisor classes are exactly the
  tetrahedron monomials $e^{(1)}_{S}e^{(2)}_{T\setminus S}$. (B) The
  pull-backs of $t_{\pm}$ along the maps $(x,y)\mapsto ax+by$ have five
  homogeneous components each, spanning $10$ dimensions; with the products
  of divisor classes they span $110$ of the $132$, with the $\Sp$-invariant
  products $16$; pull-backs along forty random pairs of diagonal matrices
  constant on each tetrahedron add nothing. (C) The classes
  $\pi_{12}-\pi_{13}$, $\pi_{12}-\pi_{23}$ and four random $\omega_{a}$
  with $a_{0}=0$ and $a_{1},a_{2},a_{3}$ not all equal lie outside the
  $110$-dimensional span, $\pi_{0}$ and $\pi_{12}+\pi_{13}+\pi_{23}$ inside
  it, and adding the four $\pi$ raises the dimension by two. (D) Pull-backs
  along forty random pairs with $\varphi_{w}=\varphi_{-w}$ and
  $\psi_{w}=\psi_{-w}$ span, with the products of divisor classes, all
  $132$ classes. (E) With $s_{t}$ the exchange of the $t$-th tensor factors,
  $s_{1}s_{2}s_{3}=-1$ on ${\textstyle\bigwedge^{2}}V$, and
  $\pi_{0},\pi_{12},\pi_{13},\pi_{23}$ are the classes of the operators
  $(1-s_{1}-s_{2}-s_{3})/4$, $(1+s_{1}+s_{2}-s_{3})/4$,
  $(1+s_{1}-s_{2}+s_{3})/4$, $(1-s_{1}+s_{2}+s_{3})/4$. (F) An $\omega_{a}$
  with $a_{0}=0$ is supported on the $12$ mixed tetrahedron monomials and on
  $24$ products of two divisor classes, $16$ of the form
  $e^{(1)}_{u}e^{(1)}_{v}e^{(2)}_{-u}e^{(2)}_{-v}$ and $8$ of the form
  $e^{(1)}_{w}e^{(1)}_{-w}e^{(2)}_{w'}e^{(2)}_{-w'}$; on each tetrahedron its
  coefficients have absolute values $|a_{2}-a_{3}|/2$, $|a_{1}-a_{3}|/2$,
  $|a_{1}-a_{2}|/2$ on the pairs of monomials indexed by $\{1,2\}$,
  $\{1,3\}$, $\{2,3\}$, checked for $\pi_{12}-\pi_{13}$ and six random
  $a$; $\pi_{0}$ and $\pi_{12}+\pi_{13}+\pi_{23}$ contain no tetrahedron
  monomial; and the component of a scalar pull-back with two factors from
  each copy has six coefficients of absolute value one. That is
  \Cref{prop:cmsource}.

\item \emph{The exact rank at a quartic field.} The computation works in the
  model of item (LV), with the basis of $H^{1}(A,\CC)$ adapted to the four
  embeddings and the $120$ basis elements of $HT^{2}$. (A) The torus that
  fixes $\theta_{1}$, $\theta_{2}$ and the four $\alpha_{\sigma}$ has
  dimension six and splits $HT^{2}$ into $34$ weight spaces, of dimensions
  one (eight), four (twenty-four, sixteen of them mixing the two real
  places) and eight (two), whose images lie in disjoint sets of monomials.
  (B) The lines have rank one. (C) On each mixed space three columns free
  of $p$ span three coordinates, and the fourth coordinates of the other
  columns are, up to sign, the sixteen $e_{ij}$ with $i,j\ge1$. (D) On each
  unmixed space of dimension four the pairs $\pm(e_{i,j},e_{i+1,j})$, or
  $\pm(e_{i,j},e_{i,j+1})$, appear after one global sign. (E) On $W_{1}$
  and $W_{2}$ the part depending on $p$ involves only the coefficients of
  one variable; the $5\times5$ minors of the $6\times8$ matrix that remains
  generate the unit ideal, its $6\times6$ minors lie in the ideals of
  $V_{+}$ and $V_{-}$, the squares of the products of generators lie in the
  ideal of the minors, and with the $2\times2$ minors of the Hankel matrix
  they generate the unit ideal; after two row operations the block is two
  unit columns, two columns reaching the new rows and the matrix $K$ of the
  proof, and the reduction of $K^{0}$ to $Q$ is confirmed at thirty-three
  points, with the ideal computations for $Q$, all by Gr\"obner bases over
  $\QQ$. (F) The formula of \Cref{thm:quarticrank}(i) agrees with the rank
  computed on all $120$ columns at forty random $p$. (G) The fifteen
  explicit $p$ of the proof have the fifteen values, and the tuples allowed
  by (C) to (E) give exactly these, the least $80$ only for constant $p$,
  and never $100$. (H) For $\omega$ with coefficients $4$, $1$, $1$, $1$ the exceptional coefficient moves as in
  \Cref{thm:quarticrank}(iv): $r=87$ at $e_{20}=2$ and $r=88$ at
  $e_{20}=3$. (I) For $\kappa_{0}=\theta_{1}+\theta_{2}$,
  $\int\alpha_{\sigma}\kappa_{0}^{6}=0$, $\int\alpha_{\sigma}\alpha_{\sigma'}
  \alpha_{\sigma''}\kappa_{0}^{2}=0$ and $\int\alpha_{\sigma}
  \alpha_{\sigma'}\kappa_{0}^{4}\ne0$ exactly when $\sigma'=\bbar\sigma$.
  (J) $\theta_{1}^{4}$ and $\theta_{2}^{4}$ are nonzero multiples of
  $\alpha_{\sigma}\alpha_{\bbar\sigma}$ for the two pairs of conjugate
  embeddings, while $\theta_{1}$, $\theta_{1}^{2}$ and $\theta_{1}^{3}$ are
  not, and $\omega+\theta_{1}^{4}/24+\theta_{2}^{4}/24$ has rank $88$, with
  $\mu=0$ and $\rho_{1}=\rho_{2}=1$, and rank $104$ when
  $\theta_{1}^{4}\theta_{2}^{4}/576$ is added.
  (K) On the $64$ elements of $H^{1}(T_{A})$ the annihilator of $\omega+p$
  consists of $F$-linear deformations and is cut out by
  $\xi\lrcorner\theta_{t}=0$ at the mid places, of dimension $8+4a$, for the
  nine patterns of exponents with random coefficients and two choices of
  $\omega$; the characters left at $88$ and a quadratic one give $8$, the
  constant and the $\theta^{4}$ shapes $16$. (L) On seventeen shapes, an
  exceptional place adds one deformation that is not $F$-linear exactly when
  $16c_{t}^{2}=u_{\sigma}u_{\bbar\sigma}$ and $\theta_{t}^{2}$ is the only
  mid monomial there, $\theta_{t}^{4}$ terms keeping it; $\theta_{1}^{4}=24\,
  \alpha_{\sigma}\alpha_{\bbar\sigma}$ and $\omega_{1}^{2}=2u_{\sigma}
  u_{\bbar\sigma}\alpha_{\sigma}\alpha_{\bbar\sigma}$; the rational
  exceptional shapes have ranks $94$ and $110$; the stabiliser of
  $\omega_{1}+c\,\theta_{1}^{2}$ in $\mathfrak{gl}(V_{\sigma}\oplus
  V_{\bbar\sigma})$ has dimension $21$, an orbit of dimension $5$ and trace
  form of signature $(12,9)$ at the exceptional value, and $15$, $4$ and
  $(8,7)$ at an ordinary one; along the locus of
  $\omega+(\theta_{1}^{2}+\theta_{2}^{2})/4$ no class of degree two and
  exactly the span of $X_{1}$, $X_{2}$ in degree four stay of Hodge type; and
  $\int X_{1}\kappa_{h}^{6}$ has one sign for six positive hermitian forms $h$
  with Gaussian rational entries on the whole tangent space at the
  exceptional value, and both signs at $c=1/10$.
  That is \Cref{thm:quarticrank}, \Cref{lem:cmweiltori}(ii) for
  $m=n=2$, \Cref{cor:quarticleast}, \Cref{prop:quarticweiltori},
  \Cref{prop:quarticlocus} and \Cref{cor:quarticlocus}.

\item \emph{Delsarte fourfolds.} (A) The sums of terms $x^{6}$, chains and
  loops in six variables have $29$ shapes, the coefficient of $t^{6}$ in the
  generating function of \Cref{prop:delsarte}. (B) For each of them $A$ is
  invertible, $A\operatorname{adj}(A)=\det(A)I$, the rows of
  $\operatorname{adj}(A)$ sum to $\det(A)/6$, and after the shift by $c$ the
  exponent vectors of $F_{A}(\psi(y))$ are those of
  $(y_{0}\cdots y_{5})^{6c}\sum_{i}y_{i}^{d}$, both for $d=|\det A|$ and for
  the least $d$, which is $24$ or $30$ for six shapes and at least $120$ for
  the others. (C) The reduced Gr\"obner basis over $\QQ$ of the partial
  derivatives has a pure power of every variable among its leading
  monomials, for the $29$ sextics and for the chains and loops of lengths
  $2$ to $6$ and degrees $3$ to $8$ completed by terms $x^{m}$; so these
  hypersurfaces are smooth. (D) For the loop sextic $\det A=5^{6}-1$,
  $\operatorname{adj}(A)$ is the circulant with first row
  $(3125,-625,125,-25,5,-1)$, and the standard monomials number $1$, $426$,
  $1751$, $426$, $1$ in degrees $0$, $6$, $12$, $18$, $24$, none in degree
  $25$, and $5^{6}$ in all. (E) The identity behind the product argument
  for loops holds for $k\le8$, $m\le12$. That is \Cref{prop:delsarte}.

\item \emph{Hypersurfaces of simplex type.} (A) For the $29$ sextic shapes
  the Smith normal form of the matrix $R$ of the differences $a_{i}-a_{0}$
  gives the invariant factors of $M/M'$; $e_{A}R^{-1}$ is integral and
  $(e_{A}/p)R^{-1}$ is not, for each prime $p$ dividing $e_{A}$;
  $[M:M']=|\det A|/6$; and $e_{A}$ divides the least $d$ of
  \Cref{prop:delsarte} and equals it except for the chain of length six and
  the loop, where it is $3125$ and $2604$. It is $6$ for the Fermat sextic,
  $24$ or $30$ for six shapes and between $120$ and $3750$ for the others.
  (B) For the Klein quartic $\det A=28$, $[M:M']=e_{A}=7$, the kernel of the
  isogeny has order $7$, and exactly three of the invariant characters of
  the Fermat septic have type $(1,0)$, the genus. (C) The sum over the
  orbits of $\PP^{r+1}$ of the Euler numbers of
  \Cref{lem:simplexchart}(i) equals $((1-m)^{r+2}-1)/m+r+2$, for the $29$
  sextics, where it is $2610$, and for $168$ chains and loops in $3$ to $6$
  variables and degrees $3$ to $8$. (D) The characters of the Fermat
  variety of degree $e_{A}$ trivial on the kernel of the isogeny, the
  classes $e_{A}xA^{-1}$ modulo $e_{A}$ with $x\in M$, form a group of order
  $[M:M']$, and those with all entries nonzero and entries summing to $e_{A}$
  number $\binom{m-1}{r+1}=h^{r,0}(X_{A})$, for the $29$ sextics and for
  $90$ chains and loops in $3$ to $5$ variables and degrees $3$ to $7$.
  (E) The same count gives the geometric genus of $40$ cyclic covers
  $y^{k}=F_{A}(x)$ in five variables, of degrees $8$ and $10$. (F) Random
  polynomials of simplex type in two and three variables have no singular
  zero on the torus, and a singular control polynomial is detected. (G) The
  smooth adapted fan of the Klein quartic has $12$ rays and unimodular
  adapted cones, and the curve meets only the three divisors of the inner
  normals, in one point each, so that its Euler number is $-7+3=-4$. That is
  \Cref{lem:simplexchart}, \Cref{thm:simplextype} and
  \Cref{cor:delsarteall}.

\item \emph{Markman's candidate for a quartic field.} The local part is an
  $\Ext$ computation over $R=\QQ[x_{1},\dots,x_{4}]$, recorded with its
  output. For ten graded modules $M$ it computes a minimal free resolution
  and its projective dimension, represents the jet class of
  $\partial/\partial x_{k}$ by the chain map with components
  $-\partial_{k}(d_{i})$, and decides for every pair $(k,l)$ whether the
  composite $\varepsilon\circ(-\partial_{k}d_{1})(-\partial_{l}d_{2})$ is a
  coboundary, by a membership test in a submodule of a free module; all the
  modules are homogeneous, so the graded computation is the local one. The
  results are those listed after \Cref{lem:divisorgerm}: the pairs in
  $\{1,2,3\}$ for a line bundle and a bundle of rank two on the curve
  $x_{1}=x_{2}=x_{3}=0$, all pairs for the skyscraper and for the ideal of a
  point in a divisor, exactly $\{1\}\times\{2,3\}$ for the ideal of a curve
  in a divisor, none for a smooth divisor and for a maximal Cohen--Macaulay
  module on a node, the four mixed pairs for two planes, the pairs in
  $\{2,3,4\}$ for the germ of \cite[Example~11.2.7]{Mar25c} at a gluing point
  and four pairs for a reflexive module on a node that is not Cohen--Macaulay.
  The global part works in the exterior algebra on
  $H^{1}(X)=U_{1}\oplus U_{2}$, with $HT^{2}(X)$ of dimension $28$ acting by
  multiplication, derivation and interior product, for
  $(q,f_{1},f_{2})=(1,2,\tfrac12)$ and $(2,3,\tfrac13)$. (A) The classes
  $\alpha_{0}$ and $\beta'$ lie in $S(0,q)$, with the coefficients of
  \Cref{prop:markmanquarticclasses}(i) and $\rho=2$. (B) The ranks of
  contraction are $12$ for $\alpha_{0}$ and $20$ for $\beta'$. (C) The
  compensated classes \eqref{eq:quarticcompensated} annihilate $S(0,q)$, and
  $e^{-\theta_{t}}x_{j}$ annihilates $e^{\theta_{t}}S(0,q)$ for three values
  of $t$. (D) The annihilator of $\alpha_{0}$ has dimension $16$ and its
  bivector components fill $\bigwedge^{2}T$; that of $\beta'$ has dimension
  $8$ and its bivector components span $\pi_{1}$, $\pi_{2}$. (E) The six
  polarised $F_{0}$-linear deformations and $x_{1}$, $x_{2}$ span the
  annihilator of $\beta'$. (F) For $F_{0}=\QQ(\sqrt5)$ and $f=(3+\sqrt5)/2$,
  $\int\alpha_{0}\beta'=-28q$, $\chi(\beta',\beta')=188q$ and
  $\chi(\alpha_{0},\alpha_{0})=8q$. (G) The rank of contraction from
  $HT^{2}(X\times X)$ into $\alpha_{0}\boxtimes\beta'$ is $96$ modulo two
  primes. (H) The first factor of \cite[Example~8.2.4]{Mar25a} has character
  $\Theta-\frac d6\Theta^{3}$ and $\chi=8d$ for $d\le12$, from the characters
  of its pieces. That is \Cref{lem:freegerm}, \Cref{lem:divisorgerm},
  \Cref{thm:quarticlocal}, \Cref{prop:markmanquarticclasses} and
  \Cref{cor:markmanquarticfails}.

\item \emph{The secant plane and the lattice of line bundles.} The binomial
  moments of \Cref{lem:binomialmoments} are computed from the coefficients
  $(a,b,-ad,-bd,ad^{2})$ of $a\,u_{\Theta}+b\,v_{\Theta}$ by solving
  \eqref{eq:stirling} in exact arithmetic, and compared with the closed
  formula of \Cref{thm:lattice} for $a\in\{1,2\}$, $b\in\{1,3\}$ and
  $1\le d\le48$. For $b=3$, $a=1$ and every $d<400$ the integrality of the
  five moments is compared with the congruence $d\equiv15,23\pmod{24}$, and
  for $d<60$ the membership of $(1,3,-d,-3d,d^{2})$ in the lattice is
  decided a second time by solving against the unimodular basis
  $e^{j\Theta}$, $0\le j\le4$. The witness
  $-23[\cO_{X}]+66[\cO_{X}(t)]-54[\cO_{X}(2t)]+20[\cO_{X}(3t)]-3[\cO_{X}(4t)]$
  of \Cref{thm:secantexists} is checked to have the moments of
  $6(u_{\Theta}+v_{\Theta})$ at $d=3$, with $6$ the least common denominator
  of the moments of $u_{\Theta}+v_{\Theta}$. For $d\in\{1,3,5,7\}$ the
  defect of $m_{4}$ is recorded, the divisibility $24\mid3(d+9)(d-1)$ of
  \Cref{rem:latticedefect} is confirmed, and the moments of $\cO_{S}$ for a
  smooth support are computed and found non-integral in $m_{4}$; finally
  $e(S)\ne[S]^{2}$ is confirmed for $1\le N\le40$ at $b=3$, for
  \Cref{lem:smoothinjective}. That is \Cref{thm:lattice},
  \Cref{cor:latticeburch}, \Cref{rem:latticedefect} and
  \Cref{lem:smoothinjective}.

\item \emph{Resolutions of rank one and two.} The total Chern class
  $c(\cI_{Z}(b\Theta))$ is computed from $u_{\Theta}+b\,v_{\Theta}$ by
  Newton's identities, symbolically in $b$ and $d$ and in exact arithmetic
  for $1\le b\le12$ and $1\le d\le60$, where it is also found to be
  integral. The rank of cup product with $\Theta^{2}$ from $H^{2}$ to
  $H^{6}$ is computed in the exterior algebra of $\ZZ^{8}$ and found to be
  $28$. At $r=1$ the class $c_{1}(E)$ killing $c_{3}(G)$ is $-\frac b3\Theta$
  and $c_{4}(G)=-\frac{1}{72}(b^{2}+d)(b^{2}+9d)\Theta^{4}$, symbolically and
  on the same grid. At $r=2$ the condition $c_{4}(G)=0$ is solved for
  $c_{2}(E)$ and $\chi(E)$ is computed; the parity condition is compared
  with $d\equiv3\pmod4$ for odd $d<400$, and an independent search over
  $c_{1}(E)=a\Theta$, $c_{2}(E)=\frac m2\Theta^{2}$ with $|a|\le12$,
  $|m|\le150$ at $b=3$ finds rank two data for $d\le40$ exactly when
  $d\equiv3\pmod4$. The total Chern class of a rank three bundle $G$ with
  $\ch(G)=\ch(E)+u_{\Theta}+3v_{\Theta}$ and $c(E)=1+\Theta^{2}$ is found to
  have $c_{4}(G)$ equal to $180$, $144$, $84$ and $0$ times
  $\frac1{24}\Theta^{4}$ at $d=1,3,5,7$, and the Whitney formula is compared
  with Newton's identities on random data. That is
  \Cref{prop:lowrankburch} and \Cref{rem:lowrankburch}.

\item \emph{Convolutions of line bundles.} On $E_{0}^{2n}$, $E_{0}=\CC/\ZZ[i]$,
  the signed sum of $e^{c_{1}(L_{\zeta})}$ over the $2\cdot4^{n-1}$ bundles
  $L_{\zeta}$ is computed in the exterior algebra over $\QQ(i)$ and found to
  be $2\cdot4^{n-1}(\alpha+\bbar\alpha)$ for $n=2,3$ and $p\in\{0,1,-3\}$,
  with the identity $\sum_{z\in\mu_{4}}ze^{D(z)}=4c'$ on one surface behind
  it. The Hankel rank of three multiples is $3$, and $r=7n^{2}-2n$, $525$ and
  $3668$ diagonal classes, and the gap $468$ are confirmed. Contraction with
  the tangent vectors of the polarised family is found to vanish on a divisor
  class only on $\CC\theta$ for $n=2,3,4$. The inequality
  $n(X+Y)\le n(X)+n(Y)$ and the bound on a sum of nondegenerate forms adding
  to zero are tested on random hermitian data. A shortest-path computation
  over all steps of degree one, with the shifts free in their parities, finds
  every closed walk of positive excess at $n=3,4$ for several placements of
  the multiples, and one of excess at most $0$ at $n=2$. On binary trees
  with up to eight leaves the number of distinct nested sets of marked
  leaves is found to be at most $u-2$, with equality. At $n=3$ the pairs of
  pieces that can carry a cup product are listed, the components $4$ and
  $16$ of their graphs are found, and the cup products for random harmonic
  components modulo $1000003$ have rank $276$, so a kernel of $204$, and
  $249$ with the classes of the multiples. Every path of \Cref{lem:harmonicmodel}
  for products of length at least three acting on the diagonal is
  enumerated with the bound of \Cref{lem:degreedrop} and filtered by the
  Maurer--Cartan equation: for $p=0$, $t=(2,4,6)$ only the paths
  $M_{2}\to M_{3}\to L_{\zeta}\to M_{1}$ remain, for the two sign patterns
  with $\sigma_{1}=\sigma_{3}=-\sigma_{2}$; the pieces $L_{\zeta}$ alone
  carry none. For $t=(-2,2,4)$ only the two sign patterns with
  $\sigma_{1}=\sigma_{2}=-\sigma_{3}$ keep paths, namely
  $M_{2}\to M_{3}\to M_{1}\to L_{\zeta}$ and
  $M_{3}\to M_{1}\to L_{\zeta}\to M_{2}$ through the sixteen $L_{\zeta}$
  with $\prod\zeta=\sigma_{3}$, which order $M_{2}$ and $M_{3}$ oppositely;
  the shifts they force are checked, and the hermitian forms of
  $L_{\zeta}\otimes M_{2}^{-1}$ give targets of dimension
  $\bigl((t_{2}-p)^{2}-1\bigr)^{3}$. For $t=(-4,-2,2)$ and $t=(-6,-4,-2)$
  the survivors of all eight sign patterns are found to be the duals of
  those for $(-2,2,4)$ and $(2,4,6)$. The search cuts branches by a lower
  bound for the excess still to come, and it is found to return the same
  paths as the exhaustive search at $t=(2,4,6)$. At $n=4$ paths of length
  three and four among the $L_{\zeta}$ are found to pass the same tests.
  That is \Cref{prop:weilpieces}, \Cref{lem:indexsum},
  \Cref{thm:fewpieces}, \Cref{lem:degreedrop}, \Cref{prop:nothingenters},
  \Cref{prop:cupkernel}, \Cref{thm:esixdiagonal},
  \Cref{prop:mixedplacement} and \Cref{rem:convolutionsopen}.

\item \emph{The fourfold products on $E_0^6$, an explicit convolution,
  and the classes no product removes.} Theta functions on $E_{0}$ of degrees
  $2$, $6$ and $14$ in the unitary gauge are checked against their automorphy
  factors, $\nabla_{\bar z}$ is found to kill them and $\nabla_{z}$ to agree
  with a finite difference, and the levels $0$, $1$, $2$ are found
  orthonormal with squared norms $k!(\pi d)^{k}$; the levels $0$ to $12$ of a
  product $s\bar\omega$ carry its norm to $10^{-8}$. On the four covers
  $\phi_{\zeta_{j}}$ the invariant sections of the five bundles of a path
  have the dimensions $|\det|$ ($1$, $35$, $3$, $24$, $9$ at $t=(2,5,6)$),
  and the theta basis of $B_{j}$ pulls back exactly into them. The selection
  rules of the levels hold, two quadrature grids agree to $10^{-15}$, and of
  the five trees with four leaves only $(x_{1}x_{2})(y\,x_{3})$ survives the
  degree count, in two assignments. In the order of
  \Cref{prop:explicitconvolution} the components and the two equations of
  the Maurer--Cartan equation are found, and the classes of degrees $1$, $2$,
  $3$ of the endomorphism complex number $908979$, $233213$ and $12142$, of
  which $16\cdot24^{3}=221184$ in degree two lie between $M_{2}$ and the
  $L_{\zeta}$. At $p=0$, $t=(2,5,6)$ the fourfold products on the $192$
  classes of type $(1,1,0)$, evaluated by quadrature in double precision,
  have numerical rank $12$ at each of the sixteen $L_{\zeta}$ and $192$
  together. (H) The same products are then evaluated in ball arithmetic at
  $128$ bits. Every integral is reduced to coefficients of holomorphic
  sections in a theta basis by the decomposition
  $(\nabla_{z}a)\,b=\bigl(d_{a}\nabla_{z}(ab)+H\bigr)/(d_{a}+d_{b})$ with
  $H$ holomorphic; the coefficients are found by interpolation at rational
  points, with the tails of the theta series bounded explicitly; the
  invariant sections on $B_{j}$ are the image of the projector of the two
  half periods, which are checked in rational arithmetic to act by commuting
  involutions, and have dimension $|\det|$ on all twelve surfaces; and the
  theta basis of $B_{j}$ is pulled back by interpolation. The values agree
  with the quadrature to $10^{-14}$. The Gram determinants of the $12$
  columns at each $L_{\zeta}$ and of all $192$ columns are intervals that
  exclude $0$, so the rank is exactly $192$ (smallest singular value $0.014$
  of the largest, radii below $10^{-25}$); the same holds with
  $x_{1}x_{2}^{\zeta}=x_{2}^{\zeta}x_{3}^{\zeta}=0$, the kernels of dimensions
  $11$ and $19$ on every surface being certified, with the relative sign of
  the two terms reversed, and at $t=(2,5,7)$. Each piece is found to
  have $3$, $6$, $7$ partners of the other parity differing in one, two,
  three coordinates, the last with $D=16$ six times and $64$ once. In the
  explicit convolution the $112$ groups $\cH^{3}$ of pieces that differ
  everywhere, $2560$ classes, are found to be reached by no element of
  degree one and to have no target of degree three, and the same is found
  for every such pair with adjacent shifts in $36$ random orders and shifts
  of the three placements of \Cref{lem:allthree}. That is
  \Cref{prop:explicitconvolution}, \Cref{lem:allthree},
  \Cref{thm:noconvolution} and \Cref{rem:convolutionsopen}; the ranks of (G)
  and (H) are not used in the paper.

\item \emph{Diagonal complete intersections of Vandermonde type.} The
  Lagrange identity $\sum_{i}w_{i}\lambda_{i}^{k}=0$ for $k\le N-1$ is checked
  in exact rational arithmetic for random distinct $\lambda_{i}$, and the
  space cut out by the $c$ equations is found to be the space of values of
  the polynomials of degree at most $r$. Every $c\times c$ minor of $957$
  matrices of Vandermonde type is found to be nonzero, and $40$ random pairs
  of diagonal equations with nonzero $2\times2$ minors are brought to the
  Vandermonde form exactly. In double precision, random points of $C^{r}$
  are found to map into $X$, where the Jacobian has rank $c$ (residual
  $10^{-14}$ at $81$ points); random points of $X$ are lifted through the
  roots of $F$, and the fibres of $\Phi$ are found to have exactly $|G|$
  points ($16$, $32$, $54$, $1536$ and $162$ for
  $(d,N,r)=(2,3,2),(2,4,2),(3,3,2),(2,4,3),(3,4,2)$). These two parts test
  a map whose properties are proved by hand, and no statement rests on
  them. The genus of $C$ is found three ways: by the Riemann--Hurwitz formula,
  by adjunction, and as the number of holomorphic forms
  $t^{j}\prod_{i}(t-\lambda_{i})^{-b_{i}/d}\,dt$ summed over the characters.
  The Euler number of $X$ from the complete intersection formula is found
  equal to the orbifold Euler number of $C^{r}/G$ in $23$ cases, and the
  middle Hodge numbers from Hirzebruch's formula for complete intersections
  are found equal to the sums over the characters of
  $\binom{\dim V_{a}^{1,0}}{p}\binom{\dim V_{a}^{0,1}}{r-p}$ for $43$ triples
  $(d,N,r)$, with $h^{2,2}=267$, $2584$, $48588$ for the fourfolds of two
  diagonal hypersurfaces of degrees $3$, $4$, $6$ in $\PP^{6}$. Over all
  characters, $\dim V_{a}$ is found to be the number of nonzero coordinates
  less two, and $\dim B_{[a]}$ to be at most five for $d\in\{3,4,6\}$,
  $N\le6$ and for $d=2$, $N\le12$, and to reach six at $N=7$ and $N=13$. The balanced orbits of
  \Cref{cor:twodiagonal}(iii) are counted: $70$, $490$ and $6125$. That is
  \Cref{thm:vandermonde}, \Cref{cor:twodiagonal} and
  \Cref{rem:vandermonde}.

\item \emph{Convolutions of the Weil pieces on $E_0^8$ at two and three
  shifts.} Each
  of the $128$ pieces $L_{\zeta}$ at $n=4$ is found to have $4$, $12$, $28$
  and $20$ partners of the other parity differing in one, two, three and
  four coordinates; the $28$ have $D=16$ twenty-four times and $D=64$ four
  times, $640$ in all, and $64\cdot640=40960$. From the hermitian forms of
  the quotients, no $H^{0}$ is found between distinct pieces, and $H^{1}$
  only between pieces of opposite parity differing in one coordinate. All
  $1020$ runs from a piece through the multiples back to a piece, distinct
  before and after one marked step, are found to lower the shift by at least
  four in the four placements $t=(2,5,6)$, $(-2,2,5)$, $(-5,-2,2)$,
  $(-6,-5,-2)$ at $p=0$. For each placement and both signs $\epsilon$, every
  chain of components that the degrees allow is enumerated at each of the
  $1792$ pairs of pieces of \Cref{thm:efourtwolevels}(ii), with the shifts
  of the multiples left free: no element of degree one reaches the
  $40960$ classes, and no term of $d_{E}$ leaves them; the identity
  $3+d_{c}-d_{e}=4-U+7D$ is checked on every chain. When one piece is moved
  to a third shift, terms that the degrees allow are found on $6$ of its
  $28$ groups, so the vanishing of $d_{E}$ needs the two shifts. With
  three consecutive shifts, every piece is found to have $18$, $24$ and
  $21$ partners of its own parity differing in two, three and four
  coordinates, with groups $H^{5}$ of dimensions adding up to $576$, $384$
  and $0$, so $960$ in all. With the pieces replaced by their three types
  and the multiples' shifts left free, every chain the degrees allow is
  enumerated in the four placements, for every split of the pieces between
  the highest and the lowest shift: the only terms of $d_{E}$ on the
  $40960$ classes land in groups $H^{5}$ between those two shifts. Piece by
  piece, for every placement with a random split and, in one placement, the
  two extreme splits, the $1792$ pairs, the absence of incoming terms and
  the targets are found again, of total dimension at most $18432$. The
  character sums of the Cayley graph of the targets are real, with least
  value $-192$, so a cut weighs at most $18432$, and the split by
  $\zeta_{3}\zeta_{4}\in\{1,i\}$ attains it. That is \Cref{lem:efourshifts},
  \Cref{thm:efourtwolevels}, \Cref{thm:efourthreelevels} and
  \Cref{rem:convolutionsopen}.
\item \emph{Very general diagonal complete intersections.} For a chain of
  $2g$ vanishing cycles with $g\le7$, the logarithms of the squared
  transvections are found to preserve the intersection form, to square to
  zero and to generate a Lie algebra of dimension $g(2g+1)$; the chains of
  $2g+1$ cycles have intersection matrices of rank $2g$ for $g\le12$. The
  invariants of $\mathfrak{sp}_{2g}$ in $\bigwedge^{r}$ are a line for $r$
  even and zero for $r$ odd ($g\le4$), and those of $\mathfrak{sl}_{n}$ vanish
  for $0<r<n$ ($n\le6$). For triple covers, the signature $(p,q)$ of an
  eigenspace and the numbers $n_{1}=2p-q+1$, $n_{2}=2q-p+1$ of branch points
  of exponents $1$ and $2$ are checked to agree with the branch data of
  Achter and Pries. On $24$ random hermitian lattices over $\ZZ[\zeta_{3}]$
  the trace form is found to have determinant
  $3^{g}\Nm(c)^{g}(\det M)^{2}$ in absolute value. The dimensions of the
  spaces of Hodge classes of the very general member are enumerated over the
  characters and agree with the closed formulas, never exceed
  $h^{r/2,r/2}$, and equal it for quadrics, for the Fermat cubics ($7$, $21$,
  $71$) and for two quadrics ($N+2$). That is \Cref{lem:vgmonodromy},
  \Cref{thm:vgvandermonde} and \Cref{rem:vandermonde}.
\item \emph{Monodromy of cyclic covers of degree three, four and six.} In
  exact arithmetic over $\QQ(\zeta_{m})$, $m\in\{3,4,6\}$, the pure braid
  group acts on the model $E_{a}$ by Fox Jacobians: each generator $A_{st}$
  acts as a pseudo-reflection of determinant $(t_{s}t_{t})^{\pm1}$ in the
  $40$ cases with five or six points, and for five to seven points the
  invariant hermitian form of the $70$ cases is unique up to a scalar and
  has signature $\{p,q\}$. In the $38$ cases with $n=3,4$ and $p,q\ge1$ the
  logarithms of the unipotent twists, closed under brackets and conjugation
  by the $A_{st}$, span $\mathfrak{sl}_{n}$. For five to nine points the
  cycle of a disc and the cycles outside it split $E_{a}$, the doubled
  braids act on the cycles outside as on $E_{a'}$ and on the disc by a root
  of unity, and the twist $A_{23}$ moves both ($61$ cases). The merge lemma
  is checked for all $57213$ admissible multisets with $7\le k\le40$
  ($k\le30$ for $m=6$), the determinant of the trace form on random
  hermitian lattices over $\ZZ[i]$, and the counts of Hodge classes of the
  very general member of degree $4$ and $6$ against the closed formula and
  against $h^{r/2,r/2}$ ($142$, $1108$ and $1752$ for the Fermat quartic
  fourfold and sixfold and the Fermat sextic fourfold). That is
  \Cref{lem:foxmodel}, \Cref{lem:cabling}, \Cref{lem:merge},
  \Cref{prop:cyclicmonodromy}, \Cref{lem:newdisc} and
  \Cref{thm:vgvandermonde}; the paper proves the base of the induction by
  hand (\Cref{lem:mergefive} and the transvection of two points), so the
  $38$ cases with $n=3,4$ are not used.
\item \emph{Groups of spread two and three on $E_0^8$, and five consecutive
  shifts.} Over the ratios of two pieces, every term of the differential with
  leaves between pieces $L_{\zeta}$ that the degrees on the four surfaces and
  \Cref{lem:degreedrop} allow is enumerated: the $13$ groups $H^{4}$ of
  spread two whose ratio moves every coordinate, one of them by $-1$, are
  reached by no coboundary and left by no term, while the $8$ groups with all
  entries $\pm i$ are reached and left; no term leaves any of the $76$ groups
  of spread three; every term leaving a group $H^{3}$ of three moved
  coordinates at spread one has chains of total drop one or two. Every run
  through the multiples, in the four placements, gives a degree
  $3+\delta+X-U+7D$ that is never the degree of its target, for $\delta=1,2,3$
  and every $X\ge0$. All $129$ subgroups $H$ of the $64$ ratios of one parity
  satisfy $|H|\,w(S\setminus H)\ge1024$, with equality only for $H=0$, and
  the algorithm of Stoer and Wagner finds the least cut, $1024$, at one
  vertex. For $40$ random splits of one parity between two values two apart,
  or three with the middle one occupied, the groups $H^{4}$ weigh at least
  $1024$. Piece by piece, with the degrees of item (LXXIV) and the shifts of
  the multiples left free, the arrangement of single shifts three apart is
  checked in two placements ($2816$ groups, $68608$ classes, left by no term
  and reached only through the loops at their two pieces), and four
  arrangements with a gap of four, whose groups $H^{3}$ beside the gap are
  reached by nothing and left by nothing. Every arrangement of the two
  parities within five consecutive values is excluded by one of the four
  results, and within six values exactly two are left up to the exchange of
  the parities. For general components on every pair, the cup products on
  the diagonal at $n=4$ have rank $3184$ modulo $1000003$ on the $3584$
  classes of the $L_{\zeta}$, with graphs $\Gamma_{\tau}$ of $4$ and $16$
  components, so their kernel is $400+84=484$. That is
  \Cref{lem:efourspread}, \Cref{lem:cayleycut}, \Cref{thm:efourspreadtwo},
  \Cref{thm:efourspreadthree}, \Cref{prop:efourgap}, \Cref{cor:efourfive},
  \Cref{prop:efourcupkernel} and \Cref{rem:convolutionsopen}.
\item \emph{The diagonal at a lonely shift, interleaved shifts, and six
  consecutive shifts on $E_0^8$.} A run from a piece through one, two or
  three distinct multiples to a piece lowers the shift by $4$, $5$, $6$,
  $12$, $13$ or $20$ in the four placements, never by $7$. With the parities
  at single shifts $s$ and $s-g$, for $g=5,7,9,11,13$, both parities on top
  and the four placements, every term on a diagonal class that the degrees
  allow is enumerated with the shifts of the multiples left free: at each
  shift there are exactly four, each runs through the three multiples and no
  other piece along the paths of the footnote to \Cref{lem:efourlonely}, the
  shifts of the piece and the multiples are pairwise incongruent modulo $4$,
  and the multiples never serve both shifts ($320$ patterns). The $28$ ratios
  that move three coordinates and change the parity weigh $640$ and generate
  the $128$ ratios; over all $636$ subgroups $H$,
  $|H|\,w(S_{3}\setminus H)\ge640$ with equality only at $H=0$, and the
  algorithm of Stoer and Wagner finds the least cut, $640$, at one vertex.
  For two random splits of the parities between $\{s,s-4\}$ and
  $\{s-1,s-5\}$, in two placements, one for each parity on top, the groups
  $H^{3}$ between adjacent shifts are reached by nothing and left by
  nothing, piece by piece with the degrees of item (LXXIV) and the shifts of
  the multiples left free, and weigh at least $640$. Every arrangement of
  the two parities within six consecutive values is excluded by the results
  of item (LXXVII), by single shifts at an odd distance or by the
  interleaved shifts, and so are single shifts at every odd distance up to
  $41$. That is \Cref{lem:efourlonely}, \Cref{prop:efoursingle},
  \Cref{prop:efourpairs}, \Cref{cor:efoursix} and
  \Cref{rem:convolutionsopen}.
\end{enumerate}

No theorem of the paper depends on these items. Items (IV) and (V) support
statements that no obstruction exists; a failure there would strengthen the
paper rather than weaken it.
The seventy-five items other than (XXIX), (XXXI) and (XL) perform $2218$
checks, all of which pass; those three, and the local part of (LXVIII), are
$\Ext$ computations over a polynomial ring, recorded with their output.
